%% Generated by Sphinx.
\def\sphinxdocclass{report}
\documentclass[a4paper,12pt,polish,openany,oneside]{sphinxmanual}
\ifdefined\pdfpxdimen
   \let\sphinxpxdimen\pdfpxdimen\else\newdimen\sphinxpxdimen
\fi \sphinxpxdimen=.75bp\relax
\ifdefined\pdfimageresolution
    \pdfimageresolution= \numexpr \dimexpr1in\relax/\sphinxpxdimen\relax
\fi
\newdimen\sphinxremdimen\sphinxremdimen = 12pt
%% let collapsible pdf bookmarks panel have high depth per default
\PassOptionsToPackage{bookmarksdepth=5}{hyperref}

\PassOptionsToPackage{booktabs}{sphinx}
\PassOptionsToPackage{colorrows}{sphinx}

\PassOptionsToPackage{warn}{textcomp}
\usepackage[utf8]{inputenc}
\ifdefined\DeclareUnicodeCharacter
% support both utf8 and utf8x syntaxes
  \ifdefined\DeclareUnicodeCharacterAsOptional
    \def\sphinxDUC#1{\DeclareUnicodeCharacter{"#1}}
  \else
    \let\sphinxDUC\DeclareUnicodeCharacter
  \fi
  \sphinxDUC{00A0}{\nobreakspace}
  \sphinxDUC{2500}{\sphinxunichar{2500}}
  \sphinxDUC{2502}{\sphinxunichar{2502}}
  \sphinxDUC{2514}{\sphinxunichar{2514}}
  \sphinxDUC{251C}{\sphinxunichar{251C}}
  \sphinxDUC{2572}{\textbackslash}
\fi
\usepackage{cmap}
\usepackage[T1]{fontenc}
\usepackage{amsmath,amssymb,amstext}
\usepackage{babel}



\usepackage{tgtermes}
\usepackage{tgheros}
\renewcommand{\ttdefault}{txtt}



\usepackage[Sonny]{fncychap}
\ChNameVar{\Large\normalfont\sffamily}
\ChTitleVar{\Large\normalfont\sffamily}
\usepackage{sphinx}

\fvset{fontsize=auto}
\usepackage{geometry}


% Include hyperref last.
\usepackage{hyperref}
% Fix anchor placement for figures with captions.
\usepackage{hypcap}% it must be loaded after hyperref.
% Set up styles of URL: it should be placed after hyperref.
\urlstyle{same}

\addto\captionspolish{\renewcommand{\contentsname}{Spis treści:}}

\usepackage{sphinxmessages}
\setcounter{tocdepth}{1}


\usepackage[utf8]{inputenc}
\usepackage[T1]{fontenc}


\title{Sprawozdanie bazy danych}
\date{28 cze 2026}
\release{}
\author{Mateusz}
\newcommand{\sphinxlogo}{\vbox{}}
\renewcommand{\releasename}{}
\makeindex
\begin{document}

\ifdefined\shorthandoff
  \ifnum\catcode`\=\string=\active\shorthandoff{=}\fi
  \ifnum\catcode`\"=\active\shorthandoff{"}\fi
\fi

\pagestyle{empty}
\sphinxmaketitle
\pagestyle{plain}
\sphinxtableofcontents
\pagestyle{normal}
\phantomsection\label{\detokenize{index::doc}}


\sphinxstepscope


\chapter{Rozdzial 1}
\label{\detokenize{rozdzial_1/index:rozdzial-1}}\label{\detokenize{rozdzial_1/index::doc}}
\sphinxAtStartPar
Wstep…

\sphinxstepscope


\chapter{Badania literaturowe}
\label{\detokenize{rozdzial_2/index:badania-literaturowe}}\label{\detokenize{rozdzial_2/index::doc}}
\sphinxstepscope


\section{Sprzęt dla bazy danych}
\label{\detokenize{rozdzial_2/rozdzial_21/index:sprzet-dla-bazy-danych}}\label{\detokenize{rozdzial_2/rozdzial_21/index::doc}}\begin{quote}\begin{description}
\sphinxlineitem{Autorzy}\begin{enumerate}
\sphinxsetlistlabels{\arabic}{enumi}{enumii}{}{.}%
\item {} 
\sphinxAtStartPar
Wiktor Głogowski

\item {} 
\sphinxAtStartPar
Olga Grześkowiak

\item {} 
\sphinxAtStartPar
Kamil Karaś

\end{enumerate}

\end{description}\end{quote}


\subsection{Wstęp}
\label{\detokenize{rozdzial_2/rozdzial_21/index:wstep}}
\sphinxAtStartPar
Współczesne systemy zarządzania bazami danych (DBMS) stanowią fundament operacyjny większości przedsiębiorstw. Efektywność ich działania zależy bezpośrednio od synergii pomiędzy warstwą oprogramowania a architekturą sprzętową, na której są one uruchamiane. Bazy danych wykazują unikalne, skrajnie wysokie wymagania względem podsystemów wejścia/wyjścia (I/O), procesora, pamięci operacyjnej oraz niezawodności zasilania.

\sphinxAtStartPar
W zależności od profilu obciążenia systemu, infrastrukturę sprzętową klasyfikuje się pod dwa główne typy zadań:
\begin{itemize}
\item {} 
\sphinxAtStartPar
\sphinxstylestrong{OLTP (Online Transaction Processing):} Systemy transakcyjne charakteryzujące się ogromną liczbą jednoczesnych, krótkich operacji zapisu i odczytu (np. systemy bankowe, sklepy internetowe). Kluczowym parametrem jest tu minimalne opóźnienie oraz wysoka liczba operacji wejścia/wyjścia na sekundę (IOPS).

\item {} 
\sphinxAtStartPar
\sphinxstylestrong{OLAP (Online Analytical Processing):} Hurtownie danych i systemy analityczne realizujące długie, złożone zapytania, które agregują miliony rekordów. Tutaj priorytetem jest przepustowość sekwencyjna oraz masowa wielordzeniowość procesorów.

\end{itemize}

\sphinxAtStartPar
Poniższy rozdział opisuje architekturę fizyczną bazy danych oraz wytyczne w sprawie doboru komponentów.


\subsection{Procesor (CPU)}
\label{\detokenize{rozdzial_2/rozdzial_21/index:procesor-cpu}}
\sphinxAtStartPar
Wybór procesora dla serwera bazodanowego determinowany jest nie tylko surową wydajnością obliczeniową, ale również architekturą pamięci podręcznej, topologią połączeń międzygniazdowych oraz modelami licencjonowania oprogramowania.


\subsubsection{Wydajność Jednowątkowa vs Wielordzeniowość}
\label{\detokenize{rozdzial_2/rozdzial_21/index:wydajnosc-jednowatkowa-vs-wielordzeniowosc}}
\sphinxAtStartPar
Decyzja o wyborze modelu procesora zależy od struktury zapytań generowanych przez aplikację:
\begin{enumerate}
\sphinxsetlistlabels{\arabic}{enumi}{enumii}{}{.}%
\item {} 
\sphinxAtStartPar
\sphinxstylestrong{Taktowanie Zegara:} W środowiskach OLTP, gdzie transakcje często wykonują się sekwencyjnie lub są blokowane przez zamki i zatrzaski, wyższe taktowanie rdzenia (np. 3.6 GHz \sphinxhyphen{} 4.0 GHz) przynosi lepsze rezultaty niż ekstremalna liczba rdzeni o niskim taktowaniu. Szybka realizacja pojedynczego wątku skraca czas trwania blokad.

\item {} 
\sphinxAtStartPar
\sphinxstylestrong{Skalowanie Wielordzeniowe:} W systemach OLAP silniki baz danych potrafią dzielić jedno zapytanie na wiele wątków. W tym przypadku procesory posiadające 64, 96 lub więcej rdzeni fizycznych (np. AMD EPYC, Intel Xeon Scalable) pozwalają na równoległe skanowanie indeksów i tabel, drastycznie przyspieszając operacje analityczne.

\end{enumerate}

\sphinxAtStartPar
\sphinxstylestrong{Wpływ Licencjonowania na Wybór CPU:} Wiodący producenci oprogramowania bazodanowego (np. Microsoft SQL Server Enterprise, Oracle) stosują model licencjonowania per\sphinxhyphen{}core. Koszt licencji na jeden rdzeń wielokrotnie przewyższa koszt zakupu samego sprzętu. Z tego powodu ekonomicznie i wydajnościowo uzasadniony jest wybór procesorów o mniejszej liczbie rdzeni, ale o maksymalnym dostępnym taktowaniu bazowym i zaawansowanej architekturze cache.


\subsubsection{Topologia NUMA (Non\sphinxhyphen{}Uniform Memory Access)}
\label{\detokenize{rozdzial_2/rozdzial_21/index:topologia-numa-non-uniform-memory-access}}
\sphinxAtStartPar
Wieloprocesorowe serwery (np. architektury dwu\sphinxhyphen{} lub czterogniazdowe) dzielą zasoby na tzw. węzły NUMA. Każdy procesor posiada swój zintegrowany kontroler pamięci i bezpośredni dostęp do fizycznie najbliższych banków RAM (pamięć lokalna).
\begin{itemize}
\item {} 
\sphinxAtStartPar
\sphinxstylestrong{Dostęp Lokalny:} Charakteryzuje się minimalnymi opóźnieniami i maksymalnym pasmem przenoszenia.

\item {} 
\sphinxAtStartPar
\sphinxstylestrong{Dostęp Zdalny:} Jeśli procesor osadzony w gnieździe 0 potrzebuje danych znajdujących się w banku pamięci przypisanym do gniazda 1, transakcja musi odbyć się za pośrednictwem magistrali międzyprocesorowej (Intel UPI lub AMD Infinity Fabric). Generuje to narzut opóźnienia zwany \sphinxstyleemphasis{NUMA penalty}.

\end{itemize}

\sphinxAtStartPar
Niewłaściwa alokacja wątków przez system operacyjny może prowadzić do ciągłego przemieszczania danych między węzłami (\sphinxstyleemphasis{NUMA bouncing}), co drastycznie obniża wydajność bazy danych. Wymagane jest stosowanie silników bazodanowych w pełni wspierających i optymalizujących architekturę NUMA (\sphinxstyleemphasis{NUMA\sphinxhyphen{}aware}).


\subsubsection{Znaczenie Pamięci Podręcznej (Cache L3)}
\label{\detokenize{rozdzial_2/rozdzial_21/index:znaczenie-pamieci-podrecznej-cache-l3}}
\sphinxAtStartPar
Struktury danych w relacyjnych bazach (drzewa B+, tabele stron) sprawiają, że procesor wykonuje ogromną liczbę operacji skoków i wyszukiwań wskaźników. Powoduje to częste chybienia pamięci podręcznej, zmuszając procesor do czekania na dane z RAM\sphinxhyphen{}u. Procesory wyposażone w powiększoną pamięć podręczną poziomu trzeciego (L3), takie jak układy z technologią 3D V\sphinxhyphen{}Cache, pozwalają na zatrzymanie krytycznych struktur indeksowych wewnątrz struktur procesora, generując ogromne zyski wydajnościowe.


\subsection{Pamięć RAM}
\label{\detokenize{rozdzial_2/rozdzial_21/index:pamiec-ram}}
\sphinxAtStartPar
Pamięć operacyjna RAM decyduje o tym, jak duża część bazy danych może być przetwarzana bezpośrednio w super\sphinxhyphen{}szybkiej pamięci półprzewodnikowej, bez konieczności odwoływania się do podsystemu dyskowego.


\subsubsection{Mechanizm Buforowania Danych}
\label{\detokenize{rozdzial_2/rozdzial_21/index:mechanizm-buforowania-danych}}
\sphinxAtStartPar
Podstawową jednostką wymiany danych między dyskiem a pamięcią w bazach danych jest strona (zazwyczaj o rozmiarze 8 KB). Silnik bazy danych alokuje obszar pamięci RAM nazywany \sphinxstylestrong{Buffer Pool} (w MS SQL Server) lub \sphinxstylestrong{Shared Buffers} (w PostgreSQL).
\begin{itemize}
\item {} 
\sphinxAtStartPar
\sphinxstylestrong{Odczyty z Pamięci (Logical Reads):} Jeśli żądana strona danych znajduje się już w RAM, odczyt następuje natychmiastowo (rzędu nanosekund).

\item {} 
\sphinxAtStartPar
\sphinxstylestrong{Odczyty z Dysku (Physical Reads):} W przypadku braku strony w buforze, wątek bazy danych zostaje zawieszony do czasu fizycznego załadowania strony z pamięci masowej, co zwiększa czas odpowiedzi o rzędy wielkości.

\end{itemize}

\sphinxAtStartPar
Docelowym stanem konfiguracyjnym dla systemów produkcyjnych jest posiadanie takiej ilości pamięci RAM, która pozwala na bezproblemowe przetrzymywanie w całości tzw. hot data set, czyli najczęściej modyfikowanych i odpytywanych danych wraz z ich indeksami.


\subsubsection{Wymóg Technologii ECC (Error\sphinxhyphen{}Correcting Code)}
\label{\detokenize{rozdzial_2/rozdzial_21/index:wymog-technologii-ecc-error-correcting-code}}
\sphinxAtStartPar
W serwerach bazodanowych niedopuszczalne jest stosowanie pamięci bez obsługi korekcji błędów (non\sphinxhyphen{}ECC). Pamięci \sphinxstylestrong{ECC} posiadają dodatkowe układy scalone pozwalające na wykrywanie i automatyczną korekcję błędów jednobitowych oraz wykrywanie błędów wielobitowych.

\sphinxAtStartPar
Wystąpienie zjawiska \sphinxstyleemphasis{bit\sphinxhyphen{}flip} (samorzutnej zmiany wartości bitu w pamięci RAM pod wpływem np. tła radiacyjnego lub zakłóceń elektromagnetycznych) w systemie bez ECC niesie katastrofalne skutki dla baz danych:
\begin{itemize}
\item {} 
\sphinxAtStartPar
\sphinxstylestrong{Cicha korupcja danych:} Błędna wartość może zostać zapisana z bufora pamięci wprost na dysk, powodując trwałe i niezauważalne zniszczenie struktury tabeli lub błędne zapisy księgowe.

\item {} 
\sphinxAtStartPar
\sphinxstylestrong{Uszkodzenie stron indeksowych:} Powoduje błędy integralności i uniemożliwia poprawne wykonywanie zapytań SQL.

\item {} 
\sphinxAtStartPar
\sphinxstylestrong{Awaria krytyczna:} W przypadku naruszenia pamięci jądra systemu operacyjnego lub krytycznego procesu bazy danych dochodzi do natychmiastowego zatrzymania serwera (Kernel Panic / BSOD).

\end{itemize}


\subsubsection{Parametry Wydajnościowe: Przepustowość i Opóźnienia}
\label{\detokenize{rozdzial_2/rozdzial_21/index:parametry-wydajnosciowe-przepustowosc-i-opoznienia}}
\sphinxAtStartPar
Współczesne standardy serwerowe opierają się na pamięciach DDR5, oferujących zaawansowaną architekturę wielokanałową (np. 8 lub 12 kanałów pamięci na jedno gniazdo CPU). Wysoka przepustowość (wyrażana w MT/s) przyspiesza operacje sortowania w pamięci, tworzenie tabel tymczasowych oraz operacje złączania typu \sphinxstyleemphasis{Hash Join}, które są na porządku dziennym w systemach analitycznych.


\subsection{Pamięć Masowa}
\label{\detokenize{rozdzial_2/rozdzial_21/index:pamiec-masowa}}
\sphinxAtStartPar
Podsystem pamięci masowej bezpośrednio determinuje ostateczną trwałość danych (cecha \sphinxstyleemphasis{Durability} z paradygmatu ACID) oraz jest najczęstszym wąskim gardłem systemu informatycznego z uwagi na mechaniczne lub interfejsowe ograniczenia prędkości.


\subsubsection{Klasyfikacja Nośników i Ich Przydatność}
\label{\detokenize{rozdzial_2/rozdzial_21/index:klasyfikacja-nosnikow-i-ich-przydatnosc}}
\sphinxAtStartPar
Współczesna inżynieria bazodanowa całkowicie odchodzi od talerzowych dysków magnetycznych HDD w środowiskach produkcyjnych na rzecz pamięci półprzewodnikowych Flash.
\begin{itemize}
\item {} 
\sphinxAtStartPar
\sphinxstylestrong{Dyski HDD (SAS 10K/15K):} Ze względu na ograniczenia mechaniczne oferują zaledwie kilkaset IOPS i wysokie opóźnienia (\textgreater{}5 ms). Ich zastosowanie ogranicza się obecnie wyłącznie do składowania archiwalnych kopii zapasowych (backups).

\item {} 
\sphinxAtStartPar
\sphinxstylestrong{Dyski SSD (SATA/SAS Enterprise):} Zapewniają stabilną wydajność na poziomie do 90 000 IOPS dla odczytu, jednak ograniczeniem staje się interfejs SATA (6 Gb/s) lub SAS\sphinxhyphen{}3 (12 Gb/s).

\item {} 
\sphinxAtStartPar
\sphinxstylestrong{Nośniki NVMe (PCIe Gen4/Gen5):} Komunikują się bezpośrednio z procesorem poprzez magistralę PCI Express, eliminując narzut protokołów SCSI/ATA. Osiągają wydajność przekraczającą 1 000 000 IOPS oraz redukują opóźnienia do wartości mikrosekundowych (\textless{}50 µs), stając się bezdyskusyjnym standardem dla baz danych.

\end{itemize}


\subsubsection{Fizyczna Separacja Architektury Plików}
\label{\detokenize{rozdzial_2/rozdzial_21/index:fizyczna-separacja-architektury-plikow}}
\sphinxAtStartPar
Poprawna konfiguracja pamięci masowej wymaga rozdzielenia różnych typów aktywności wejścia/wyjścia na odseparowane fizycznie wolumeny dyskowe:
\begin{enumerate}
\sphinxsetlistlabels{\arabic}{enumi}{enumii}{}{.}%
\item {} 
\sphinxAtStartPar
\sphinxstylestrong{Dziennik Transakcji:} Każda operacja modyfikacji danych (INSERT, UPDATE, DELETE) musi zostać najpierw sekwencyjnie zapisana w dzienniku transakcji, zanim baza potwierdzi jej pomyślne wykonanie. Ruch ten ma charakter niemal wyłącznie sekwencyjnego zapisu o krytycznym znaczeniu dla opóźnienia aplikacji. Dziennik transakcji powinien znajdować się na najszybszych dedykowanych nośnikach o minimalnym opóźnieniu zapisu.

\item {} 
\sphinxAtStartPar
\sphinxstylestrong{Pliki Danych:} Przechowują właściwe tabele i indeksy. Ruch na tych plikach ma charakter wysoce losowy (zarówno odczyty, jak i zapisy wywoływane przez proces \sphinxstyleemphasis{Checkpoint}, który zrzuca zmodyfikowane strony z RAM na dysk). Wymagają macierzy zoptymalizowanej pod kątem wysokiej liczby losowych IOPS.

\item {} 
\sphinxAtStartPar
\sphinxstylestrong{Baza Tymczasowa (TempDB / Swapping Space):} Wykorzystywana przez silnik bazy do przechowywania obiektów tymczasowych, wyników pośrednich podzapytań oraz sortowania. Generuje bardzo intensywny, mieszany ruch. Umieszcza się ją na ultra\sphinxhyphen{}szybkich dyskach o wysokiej wytrzymałości.

\end{enumerate}


\subsubsection{Żywotność i Wytrzymałość Zapisu (DWPD)}
\label{\detokenize{rozdzial_2/rozdzial_21/index:zywotnosc-i-wytrzymalosc-zapisu-dwpd}}
\sphinxAtStartPar
Bazy danych nieustannie modyfikują i nadpisują sektory. Z tego powodu stosowanie dysków klasy konsumenckiej grozi ich błyskawicznym zużyciem. Przy doborze nośników kluczowym parametrem jest \sphinxstylestrong{DWPD} (\sphinxstyleemphasis{Drive Writes Per Day}) \textendash{} wskaźnik określający, ile razy dziennie można zapisać całą pojemność dysku przez okres gwarancyjny (zazwyczaj 5 lat). Do baz danych wymagane są nośniki klasy \sphinxstyleemphasis{Enterprise Mixed\sphinxhyphen{}Use} lub \sphinxstyleemphasis{Write\sphinxhyphen{}Intensive} o współczynniku DWPD wynoszącym minimum 3\sphinxhyphen{}5.


\subsection{Kontrolery i Interfejsy Pamięci}
\label{\detokenize{rozdzial_2/rozdzial_21/index:kontrolery-i-interfejsy-pamieci}}
\sphinxAtStartPar
Kontrolery realizują zadania pośrednictwa, zarządzania macierzami nadmiarowymi (RAID) oraz gwarantują integralność zapisu.


\subsubsection{Sprzętowy RAID vs Software\sphinxhyphen{}Defined Storage}
\label{\detokenize{rozdzial_2/rozdzial_21/index:sprzetowy-raid-vs-software-defined-storage}}
\sphinxAtStartPar
Przez lata standardem były sprzętowe kontrolery RAID (\sphinxstyleemphasis{Hardware RAID}) wyposażone we własny procesor i pamięć cache. Odciążają one procesor główny serwera z obliczeń sum kontrolnych (szczególnie parzystości w RAID 5/6).

\sphinxAtStartPar
W przypadku nowoczesnych pamięci NVMe sprzętowe kontrolery starszego typu stawały się wąskim gardłem, blokując bezpośredni dostęp do linii PCIe. Współcześnie odchodzi się od nich na rzecz:
\begin{itemize}
\item {} 
\sphinxAtStartPar
\sphinxstylestrong{Bezpośredniego podłączenia NVMe (Direct Attached PCIe)} i zarządzania programowego na poziomie zaawansowanych systemów plików (np. ZFS) lub warstw logicznych OS (np. Linux mdadm / LVM).

\item {} 
\sphinxAtStartPar
Modernistycznych, dedykowanych sprzętowych kontrolerów NVMe RAID klasy Enterprise, zdolnych do natywnego przetwarzania protokołu NVMe z pełną prędkością magistrali PCIe Gen5.

\end{itemize}

\sphinxAtStartPar
\sphinxstylestrong{Pamięć Podtrzymywana Bateryjnie (BBU/FBWC):}
W przypadku korzystania ze sprzętowego kontrolera RAID z włączoną pamięcią podręczną zapisu (\sphinxstyleemphasis{Write\sphinxhyphen{}Back Cache}), \sphinxstylestrong{bezwzględnym warunkiem bezpieczeństwa} jest obecność modułu \sphinxstylestrong{BBU} (\sphinxstyleemphasis{Battery Backup Unit}) lub \sphinxstylestrong{FBWC} (\sphinxstyleemphasis{Flash\sphinxhyphen{}Backed Write Cache}). W razie awarii zasilania serwera, moduł ten podtrzymuje dane w pamięci cache kontrolera (lub przepisuje je do dedykowanej pamięci flash), zapobiegając utracie nieutrwalonych transakcji i uszkodzeniu struktury macierzy.


\subsubsection{Wybór Poziomu RAID pod Kątem Baz Danych}
\label{\detokenize{rozdzial_2/rozdzial_21/index:wybor-poziomu-raid-pod-katem-baz-danych}}\begin{itemize}
\item {} 
\sphinxAtStartPar
\sphinxstylestrong{RAID 10 (Striped Mirror):} Rekomendowany układ dla większości systemów bazodanowych. Łączy zalety wydajnościowe paskowania (RAID 0) i bezpieczeństwa lustrzanego (RAID 1). Zapewnia doskonałą wydajność losowego zapisu, ponieważ nie generuje tzw. narzutu parzystości (\sphinxstyleemphasis{write penalty}).

\item {} 
\sphinxAtStartPar
\sphinxstylestrong{RAID 5 / RAID 6:} Zapewniają lepsze wykorzystanie przestrzeni dyskowej, ale każda operacja zapisu wymaga odczytu danych, obliczenia nowej parzystości i ponownego zapisu (\sphinxstyleemphasis{kara za zapis}). Są \sphinxstylestrong{wysoce niewskazane dla baz transakcyjnych OLTP}. Mogą być stosowane jedynie w hurtowniach danych (OLAP), gdzie operacje zapisu są sekwencyjne i rzadkie (procesy ETL).

\end{itemize}


\subsection{Zasilanie i Niezawodność}
\label{\detokenize{rozdzial_2/rozdzial_21/index:zasilanie-i-niezawodnosc}}
\sphinxAtStartPar
Podsystem zasilania w infrastrukturze bazodanowej traktowany jest jako integralny element ochrony integralności logicznej danych, a nie tylko komponent ciągłości działania.


\subsubsection{Nadmiarowość Zasilaczy i Niezależne Linie Zasilające}
\label{\detokenize{rozdzial_2/rozdzial_21/index:nadmiarowosc-zasilaczy-i-niezalezne-linie-zasilajace}}
\sphinxAtStartPar
Serwer bazodanowy musi posiadać zasilacze nadmiarowe typu \sphinxstyleemphasis{Hot\sphinxhyphen{}Swap} (z możliwością wymiany podczas pracy) w konfiguracji minimum \sphinxstylestrong{N+1} lub optymalnie \sphinxstylestrong{2N} (Full Redundancy).

\sphinxAtStartPar
Zasilacze muszą być podpięte do fizycznie odseparowanych źródeł zasilania:
\begin{itemize}
\item {} 
\sphinxAtStartPar
\sphinxstylestrong{Zasilacz A:} Podłączony do Linii zasilającej A oraz dedykowanego modułu UPS A.

\item {} 
\sphinxAtStartPar
\sphinxstylestrong{Zasilacz B:} Podłączony do Linii zasilającej B oraz dedykowanego modułu UPS B.

\end{itemize}

\sphinxAtStartPar
Taka topologia gwarantuje, że awaria jednego zasilacza, uszkodzenie kabla lub całkowity zanik napięcia na jednej z linii energetycznych nie spowoduje wyłączenia serwera.


\subsubsection{Infrastruktura UPS i Agregaty Prądotwórcze}
\label{\detokenize{rozdzial_2/rozdzial_21/index:infrastruktura-ups-i-agregaty-pradotworcze}}
\sphinxAtStartPar
Wszelkie systemy serwerowe DBMS muszą współpracować z zasilaczami awaryjnymi UPS działającymi w topologii \sphinxstyleemphasis{Online} (podwójne przetwarzanie). Zapewniają one idealną sinusoidę napięcia, eliminują przepięcia i gwarantują czas podtrzymania niezbędny do automatycznego uruchomienia spalinowych agregatów prądotwórczych o odpowiedniej mocy.


\subsubsection{Zapobieganie Zjawisku „Torn Writes”}
\label{\detokenize{rozdzial_2/rozdzial_21/index:zapobieganie-zjawisku-torn-writes}}
\sphinxAtStartPar
Zjawisko \sphinxstyleemphasis{Torn Write} (rozbity/częściowy zapis) występuje, gdy system operacyjny przesyła do zapisu stronę danych (np. 8 KB), a w trakcie fizycznego zapisu na komórkę pamięci Flash dochodzi do nagłego odcięcia zasilania. W efekcie na dysku ląduje np. tylko 4 KB danych, co powoduje bezpowrotne uszkodzenie struktury tej strony i błąd sumy kontrolnej przy próbie kolejnego uruchomienia bazy danych.

\sphinxAtStartPar
Poza mechanizmami programowymi (takimi jak \sphinxstyleemphasis{Doublewrite Buffer} w silniku InnoDB/MySQL), na poziomie sprzętowym stosuje się technologię \sphinxstylestrong{PLP (Power Loss Protection)}. Dyski klasy Enterprise posiadają wbudowane banki kondensatorów tantalowych. W momencie wykrycia spadku napięcia, energia z kondensatorów pozwala na podtrzymanie pracy kontrolera dysku przez czas wymagany do bezpiecznego opróżnienia pamięci podręcznej RAM dysku i utrwalenia wszystkich danych w nieulotnych komórkach pamięci NAND Flash.


\subsection{Podsumowanie}
\label{\detokenize{rozdzial_2/rozdzial_21/index:podsumowanie}}
\sphinxAtStartPar
Dobór sprzętu dla systemów zarządzania bazami danych musi być ściśle dopasowany do profilu obciążenia, z wyraźnym podziałem na transakcyjne środowiska OLTP oraz analityczne systemy OLAP. Wybór procesora stanowi kompromis pomiędzy taktowaniem a wielordzeniowością ze względu na koszty licencyjne, podczas gdy w przypadku pamięci operacyjnej absolutnym wymogiem jest zastosowanie technologii ECC chroniącej przed cichą korupcją danych. W obszarze pamięci masowej standardem stały się nośniki półprzewodnikowe NVMe o wysokiej wytrzymałości (DWPD), które ze względów wydajnościowych wymagają fizycznej separacji wolumenów dla dzienników transakcji i plików danych. Do zabezpieczenia pracy dysków rekomenduje się użycie konfiguracji RAID 10 w połączeniu z nowoczesnymi kontrolerami sprzętowymi lub rozwiązaniami programowymi. Całość infrastruktury musi być kategorycznie chroniona przed awariami poprzez nadmiarowe zasilacze w konfiguracji 2N, systemy UPS w topologii Online oraz dyski wyposażone w sprzętową ochronę PLP eliminującą zjawisko uszkodzonych zapisów.

\sphinxstepscope


\section{Konfiguracja bazy danych PostgreSQL}
\label{\detokenize{rozdzial_2/rozdzial_22/index:konfiguracja-bazy-danych-postgresql}}\label{\detokenize{rozdzial_2/rozdzial_22/index::doc}}

\subsection{Wstęp}
\label{\detokenize{rozdzial_2/rozdzial_22/index:wstep}}
\sphinxAtStartPar
Rozdział opisuje konfigurację połączenia z bazą danych \sphinxstylestrong{PostgreSQL} \textendash{} dojrzałym, open\sphinxhyphen{}source’owym systemem zarządzania relacyjnymi bazami danych (RDBMS), który wyróżnia się pełną zgodnością ze standardem SQL oraz silnym naciskiem na rozszerzalność i niezawodność.

\sphinxAtStartPar
\sphinxstylestrong{Wymagania:}
\begin{itemize}
\item {} 
\sphinxAtStartPar
dostęp do instancji PostgreSQL,

\item {} 
\sphinxAtStartPar
zmienne środowiskowe,

\item {} 
\sphinxAtStartPar
narzędzia projektowe (np. \sphinxcode{\sphinxupquote{psql}}, \sphinxcode{\sphinxupquote{pg\_isready}}).

\end{itemize}


\subsection{Podstawowe parametry}
\label{\detokenize{rozdzial_2/rozdzial_22/index:podstawowe-parametry}}
\sphinxAtStartPar
Każde połączenie z PostgreSQL wymaga:
\begin{itemize}
\item {} 
\sphinxAtStartPar
hosta i portu (domyślnie \sphinxcode{\sphinxupquote{localhost}} i \sphinxcode{\sphinxupquote{5432}}),

\item {} 
\sphinxAtStartPar
nazwy bazy danych,

\item {} 
\sphinxAtStartPar
użytkownika i hasła.

\end{itemize}

\sphinxAtStartPar
Dane wrażliwe (np. hasła) przechowuj w \sphinxstylestrong{zmiennych środowiskowych} \textendash{} to oddziela konfigurację od kodu.
Pamiętaj: zmienne środowiskowe nie są mechanizmem bezpieczeństwa; w produkcji uzupełnij je o dedykowane
zarządzanie sekretami (np. HashiCorp Vault, AWS Secrets Manager).


\subsection{Lokalizacja i struktura katalogów PostgreSQL}
\label{\detokenize{rozdzial_2/rozdzial_22/index:lokalizacja-i-struktura-katalogow-postgresql}}
\sphinxAtStartPar
Pliki konfiguracyjne PostgreSQL znajdują się w katalogu danych (\sphinxcode{\sphinxupquote{PGDATA}}).

\sphinxAtStartPar
\sphinxstylestrong{Główne pliki konfiguracyjne:}
\begin{itemize}
\item {} 
\sphinxAtStartPar
\sphinxcode{\sphinxupquote{postgresql.conf}} \textendash{} podstawowa konfiguracja serwera (port, pamięć, logi)

\item {} 
\sphinxAtStartPar
\sphinxcode{\sphinxupquote{pg\_hba.conf}} \textendash{} kontrola dostępu (kto i skąd może się łączyć)

\item {} 
\sphinxAtStartPar
\sphinxcode{\sphinxupquote{pg\_ident.conf}} \textendash{} mapowanie użytkowników systemowych na role PostgreSQL

\end{itemize}

\sphinxAtStartPar
\sphinxstylestrong{Domyślne lokalizacje:}


\begin{savenotes}\sphinxattablestart
\sphinxthistablewithglobalstyle
\centering
\sphinxcapstartof{table}
\sphinxthecaptionisattop
\sphinxcaption{Domyślne lokalizacje}\label{\detokenize{rozdzial_2/rozdzial_22/index:id6}}
\sphinxaftertopcaption
\begin{tabular}[t]{\X{30}{100}\X{70}{100}}
\sphinxtoprule
\begin{varwidth}[t]{\sphinxcolwidth{1}{2}}
\sphinxstyletheadfamily \sphinxAtStartPar
System
\sphinxbeforeendvarwidth
\end{varwidth}%
&\begin{varwidth}[t]{\sphinxcolwidth{1}{2}}
\sphinxstyletheadfamily \sphinxAtStartPar
Katalog danych
\sphinxbeforeendvarwidth
\end{varwidth}%
\\
\sphinxmidrule
\sphinxtableatstartofbodyhook\begin{varwidth}[t]{\sphinxcolwidth{1}{2}}
\sphinxAtStartPar
Linux (RPM/DEB)
\sphinxbeforeendvarwidth
\end{varwidth}%
&\begin{varwidth}[t]{\sphinxcolwidth{1}{2}}
\sphinxAtStartPar
\sphinxcode{\sphinxupquote{/var/lib/pgsql/data/}} lub \sphinxcode{\sphinxupquote{/var/lib/postgresql/*/main/}}
\sphinxbeforeendvarwidth
\end{varwidth}%
\\
\sphinxhline\begin{varwidth}[t]{\sphinxcolwidth{1}{2}}
\sphinxAtStartPar
Linux (kompilacja źródłowa)
\sphinxbeforeendvarwidth
\end{varwidth}%
&\begin{varwidth}[t]{\sphinxcolwidth{1}{2}}
\sphinxAtStartPar
\sphinxcode{\sphinxupquote{/usr/local/pgsql/data/}}
\sphinxbeforeendvarwidth
\end{varwidth}%
\\
\sphinxhline\begin{varwidth}[t]{\sphinxcolwidth{1}{2}}
\sphinxAtStartPar
Windows
\sphinxbeforeendvarwidth
\end{varwidth}%
&\begin{varwidth}[t]{\sphinxcolwidth{1}{2}}
\sphinxAtStartPar
\sphinxcode{\sphinxupquote{C:\textbackslash{}Program Files\textbackslash{}PostgreSQL\textbackslash{}\textless{}version\textgreater{}\textbackslash{}data\textbackslash{}}}
\sphinxbeforeendvarwidth
\end{varwidth}%
\\
\sphinxhline\begin{varwidth}[t]{\sphinxcolwidth{1}{2}}
\sphinxAtStartPar
Docker
\sphinxbeforeendvarwidth
\end{varwidth}%
&\begin{varwidth}[t]{\sphinxcolwidth{1}{2}}
\sphinxAtStartPar
\sphinxcode{\sphinxupquote{/var/lib/postgresql/data/}} (w kontenerze)
\sphinxbeforeendvarwidth
\end{varwidth}%
\\
\sphinxbottomrule
\end{tabular}
\sphinxtableafterendhook\par
\sphinxattableend\end{savenotes}

\sphinxAtStartPar
\sphinxstylestrong{Zmiana katalogu danych:}

\begin{sphinxVerbatim}[commandchars=\\\{\}]
\PYG{c+c1}{\PYGZsh{} Inicjalizacja nowego katalogu}
initdb\PYG{+w}{ }\PYGZhy{}D\PYG{+w}{ }/ścieżka/do/katalogu

\PYG{c+c1}{\PYGZsh{} Uruchomienie z niestandardowym PGDATA}
pg\PYGZus{}ctl\PYG{+w}{ }\PYGZhy{}D\PYG{+w}{ }/ścieżka/do/katalogu\PYG{+w}{ }start
\end{sphinxVerbatim}

\sphinxAtStartPar
\sphinxstylestrong{Struktura katalogu danych:}

\begin{sphinxVerbatim}[commandchars=\\\{\}]
\PYGZdl{}PGDATA/
├── postgresql.conf   \PYGZsh{} główny plik konfiguracyjny
├── pg\PYGZus{}hba.conf       \PYGZsh{} polityka dostępu
├── pg\PYGZus{}ident.conf     \PYGZsh{} mapowanie tożsamości
├── base/             \PYGZsh{} rzeczywiste bazy danych
├── global/           \PYGZsh{} tabele systemowe
├── log/              \PYGZsh{} logi serwera
└── pg\PYGZus{}wal/           \PYGZsh{} Write\PYGZhy{}Ahead Log (WAL)
\end{sphinxVerbatim}

\sphinxAtStartPar
\sphinxstylestrong{Sprawdzenie aktualnej lokalizacji:}

\begin{sphinxVerbatim}[commandchars=\\\{\}]
\PYG{k}{SHOW}\PYG{+w}{ }\PYG{n}{data\PYGZus{}directory}\PYG{p}{;}
\PYG{k}{SHOW}\PYG{+w}{ }\PYG{n}{config\PYGZus{}file}\PYG{p}{;}
\PYG{k}{SHOW}\PYG{+w}{ }\PYG{n}{hba\PYGZus{}file}\PYG{p}{;}
\end{sphinxVerbatim}


\subsection{Środowiska i profile}
\label{\detokenize{rozdzial_2/rozdzial_22/index:srodowiska-i-profile}}
\sphinxAtStartPar
Przełączanie środowisk PostgreSQL (dev/test/prod) realizuj przez:
\begin{itemize}
\item {} 
\sphinxAtStartPar
zmienne środowiskowe (np. \sphinxcode{\sphinxupquote{PGHOST}}, \sphinxcode{\sphinxupquote{PGPORT}}, \sphinxcode{\sphinxupquote{PGDATABASE}}, \sphinxcode{\sphinxupquote{PGUSER}}, \sphinxcode{\sphinxupquote{PGPASSWORD}}),

\item {} 
\sphinxAtStartPar
profile konfiguracyjne,

\item {} 
\sphinxAtStartPar
osobne pliki \sphinxcode{\sphinxupquote{.env}} dla każdego środowiska.

\end{itemize}


\subsection{Automatyzacja}
\label{\detokenize{rozdzial_2/rozdzial_22/index:automatyzacja}}
\sphinxAtStartPar
Zalecenia dla PostgreSQL:
\begin{itemize}
\item {} 
\sphinxAtStartPar
walidacja konfiguracji przy starcie za pomocą \sphinxcode{\sphinxupquote{pg\_isready}},

\item {} 
\sphinxAtStartPar
skrypty sprawdzające kompletność zmiennych środowiskowych,

\item {} 
\sphinxAtStartPar
integracja z CI/CD (np. GitHub Actions, GitLab CI).

\end{itemize}

\sphinxAtStartPar
Dokumentację generuj automatycznie (Sphinx).


\subsection{Przykłady}
\label{\detokenize{rozdzial_2/rozdzial_22/index:przyklady}}
\sphinxAtStartPar
\sphinxstylestrong{PostgreSQL \textendash{} zmienne środowiskowe (.env)}

\begin{sphinxVerbatim}[commandchars=\\\{\}]
PGHOST=localhost
PGPORT=5432
PGDATABASE=aplikacja
PGUSER=app\PYGZus{}user
PGPASSWORD=\PYGZdl{}\PYGZob{}DB\PYGZus{}PASSWORD\PYGZcb{}
\end{sphinxVerbatim}

\sphinxAtStartPar
\sphinxstylestrong{PostgreSQL \textendash{} URL połączenia}

\begin{sphinxVerbatim}[commandchars=\\\{\}]
DATABASE\PYGZus{}URL=postgresql://\PYGZdl{}\PYGZob{}PGUSER\PYGZcb{}:\PYGZdl{}\PYGZob{}PGPASSWORD\PYGZcb{}@\PYGZdl{}\PYGZob{}PGHOST\PYGZcb{}:\PYGZdl{}\PYGZob{}PGPORT\PYGZcb{}/\PYGZdl{}\PYGZob{}PGDATABASE\PYGZcb{}
\end{sphinxVerbatim}

\sphinxAtStartPar
\sphinxstylestrong{PostgreSQL \textendash{} sprawdzenie połączenia}

\begin{sphinxVerbatim}[commandchars=\\\{\}]
pg\PYGZus{}isready \PYGZhy{}h \PYGZdl{}\PYGZob{}PGHOST\PYGZcb{} \PYGZhy{}p \PYGZdl{}\PYGZob{}PGPORT\PYGZcb{} \PYGZhy{}U \PYGZdl{}\PYGZob{}PGUSER\PYGZcb{} \PYGZhy{}d \PYGZdl{}\PYGZob{}PGDATABASE\PYGZcb{}
\end{sphinxVerbatim}


\subsection{Najlepsze praktyki}
\label{\detokenize{rozdzial_2/rozdzial_22/index:najlepsze-praktyki}}\begin{enumerate}
\sphinxsetlistlabels{\arabic}{enumi}{enumii}{}{.}%
\item {} 
\sphinxAtStartPar
Nie przechowuj sekretów w repozytorium.

\item {} 
\sphinxAtStartPar
Stosuj \sphinxcode{\sphinxupquote{.env.example}} zamiast rzeczywistych danych.

\item {} 
\sphinxAtStartPar
Używaj standardowych zmiennych środowiskowych PostgreSQL (\sphinxcode{\sphinxupquote{PG*}}).

\item {} 
\sphinxAtStartPar
Waliduj konfigurację przed uruchomieniem (\sphinxcode{\sphinxupquote{pg\_isready}}).

\item {} 
\sphinxAtStartPar
W środowiskach produkcyjnych używaj połączeń SSL/TLS (\sphinxcode{\sphinxupquote{sslmode=require}}).

\end{enumerate}


\subsection{Podsumowanie}
\label{\detokenize{rozdzial_2/rozdzial_22/index:podsumowanie}}
\sphinxAtStartPar
Poprawna konfiguracja połączenia z PostgreSQL zwiększa bezpieczeństwo i przenośność aplikacji między środowiskami. Standardowe zmienne środowiskowe oraz narzędzia takie jak \sphinxcode{\sphinxupquote{pg\_isready}} ułatwiają automatyzację i walidację.


\subsection{Szczegółowe omówienie plików konfiguracyjnych}
\label{\detokenize{rozdzial_2/rozdzial_22/index:szczegolowe-omowienie-plikow-konfiguracyjnych}}
\sphinxAtStartPar
Po omówieniu podstawowych wymagań środowiskowych można przejść do szczegółowej konfiguracji PostgreSQL.

\sphinxAtStartPar
W kolejnych sekcjach przedstawiono najważniejsze parametry konfiguracyjne dostępne w plikach:
\begin{itemize}
\item {} 
\sphinxAtStartPar
\sphinxcode{\sphinxupquote{postgresql.conf}} \textendash{} konfiguracja działania serwera,

\item {} 
\sphinxAtStartPar
\sphinxcode{\sphinxupquote{pg\_hba.conf}} \textendash{} kontrola dostępu do bazy danych,

\item {} 
\sphinxAtStartPar
\sphinxcode{\sphinxupquote{pg\_ident.conf}} \textendash{} mapowanie użytkowników systemowych na role PostgreSQL.

\end{itemize}


\subsection{postgresql.conf}
\label{\detokenize{rozdzial_2/rozdzial_22/index:postgresql-conf}}
\sphinxAtStartPar
Plik \sphinxcode{\sphinxupquote{postgresql.conf}} jest głównym plikiem konfiguracyjnym klastra PostgreSQL.
Zawiera ustawienia wpływające na działanie serwera, wydajność, bezpieczeństwo, logowanie oraz replikację.

\sphinxAtStartPar
Lokalizację aktywnego pliku można sprawdzić poleceniem:

\begin{sphinxVerbatim}[commandchars=\\\{\}]
\PYG{k}{SHOW}\PYG{+w}{ }\PYG{n}{config\PYGZus{}file}\PYG{p}{;}
\end{sphinxVerbatim}


\subsubsection{Najważniejsze parametry}
\label{\detokenize{rozdzial_2/rozdzial_22/index:najwazniejsze-parametry}}

\paragraph{Połączenia sieciowe}
\label{\detokenize{rozdzial_2/rozdzial_22/index:polaczenia-sieciowe}}\begin{description}
\sphinxlineitem{\sphinxcode{\sphinxupquote{listen\_addresses}}}
\sphinxAtStartPar
Określa adresy IP, na których PostgreSQL nasłuchuje połączeń.

\sphinxAtStartPar
Najczęstsze wartości:
\begin{itemize}
\item {} 
\sphinxAtStartPar
\sphinxcode{\sphinxupquote{localhost}} \textendash{} tylko połączenia lokalne,

\item {} 
\sphinxAtStartPar
\sphinxcode{\sphinxupquote{*}} \textendash{} wszystkie interfejsy sieciowe,

\item {} 
\sphinxAtStartPar
lista adresów, np. \sphinxcode{\sphinxupquote{\textquotesingle{}192.168.1.10,localhost\textquotesingle{}}}.

\end{itemize}

\sphinxlineitem{\sphinxcode{\sphinxupquote{port}}}
\sphinxAtStartPar
Port TCP używany przez PostgreSQL.

\sphinxAtStartPar
Domyślna wartość:

\begin{sphinxVerbatim}[commandchars=\\\{\}]
\PYG{n+na}{port}\PYG{+w}{ }\PYG{o}{=}\PYG{+w}{ }\PYG{l+s}{5432}
\end{sphinxVerbatim}

\sphinxlineitem{\sphinxcode{\sphinxupquote{max\_connections}}}
\sphinxAtStartPar
Maksymalna liczba jednoczesnych połączeń z bazą danych.

\sphinxAtStartPar
Zbyt wysoka wartość może zwiększać zużycie pamięci.

\end{description}


\paragraph{Pamięć i wydajność}
\label{\detokenize{rozdzial_2/rozdzial_22/index:pamiec-i-wydajnosc}}\begin{description}
\sphinxlineitem{\sphinxcode{\sphinxupquote{shared\_buffers}}}
\sphinxAtStartPar
Ilość pamięci RAM używanej przez PostgreSQL do buforowania danych.

\sphinxAtStartPar
Typowe wartości:
\begin{itemize}
\item {} 
\sphinxAtStartPar
128MB \textendash{} środowiska testowe,

\item {} 
\sphinxAtStartPar
25\% pamięci RAM \textendash{} środowiska produkcyjne.

\end{itemize}

\sphinxlineitem{\sphinxcode{\sphinxupquote{work\_mem}}}
\sphinxAtStartPar
Ilość pamięci wykorzystywanej przez operacje sortowania i zapytania.

\sphinxAtStartPar
Parametr jest przydzielany dla pojedynczej operacji.

\sphinxlineitem{\sphinxcode{\sphinxupquote{maintenance\_work\_mem}}}
\sphinxAtStartPar
Pamięć używana podczas operacji administracyjnych, np.:
\begin{itemize}
\item {} 
\sphinxAtStartPar
\sphinxcode{\sphinxupquote{VACUUM}},

\item {} 
\sphinxAtStartPar
\sphinxcode{\sphinxupquote{CREATE INDEX}},

\item {} 
\sphinxAtStartPar
\sphinxcode{\sphinxupquote{ALTER TABLE}}.

\end{itemize}

\sphinxlineitem{\sphinxcode{\sphinxupquote{effective\_cache\_size}}}
\sphinxAtStartPar
Szacowany rozmiar pamięci podręcznej dostępnej dla PostgreSQL.

\sphinxAtStartPar
Parametr pomaga optymalizatorowi zapytań wybierać odpowiednie plany wykonania.

\end{description}


\paragraph{Logowanie}
\label{\detokenize{rozdzial_2/rozdzial_22/index:logowanie}}\begin{description}
\sphinxlineitem{\sphinxcode{\sphinxupquote{logging\_collector}}}
\sphinxAtStartPar
Włącza zapisywanie logów do plików.

\sphinxAtStartPar
Dostępne wartości:
\begin{itemize}
\item {} 
\sphinxAtStartPar
\sphinxcode{\sphinxupquote{on}},

\item {} 
\sphinxAtStartPar
\sphinxcode{\sphinxupquote{off}}.

\end{itemize}

\sphinxlineitem{\sphinxcode{\sphinxupquote{log\_directory}}}
\sphinxAtStartPar
Katalog przechowywania logów.

\sphinxlineitem{\sphinxcode{\sphinxupquote{log\_filename}}}
\sphinxAtStartPar
Wzorzec nazw plików logów.

\sphinxAtStartPar
Przykład:

\begin{sphinxVerbatim}[commandchars=\\\{\}]
\PYG{n+na}{log\PYGZus{}filename}\PYG{+w}{ }\PYG{o}{=}\PYG{+w}{ }\PYG{l+s}{\PYGZsq{}}\PYG{l+s}{postgresql\PYGZhy{}\PYGZpc{}Y\PYGZhy{}\PYGZpc{}m\PYGZhy{}\PYGZpc{}d.log}\PYG{l+s}{\PYGZsq{}}
\end{sphinxVerbatim}

\sphinxlineitem{\sphinxcode{\sphinxupquote{log\_min\_duration\_statement}}}
\sphinxAtStartPar
Określa minimalny czas wykonania zapytania (w ms), po którym zostanie ono zapisane w logach.

\sphinxAtStartPar
Przykład:

\begin{sphinxVerbatim}[commandchars=\\\{\}]
\PYG{n+na}{log\PYGZus{}min\PYGZus{}duration\PYGZus{}statement}\PYG{+w}{ }\PYG{o}{=}\PYG{+w}{ }\PYG{l+s}{1000}
\end{sphinxVerbatim}

\sphinxAtStartPar
Wartość \sphinxcode{\sphinxupquote{1000}} oznacza logowanie zapytań trwających dłużej niż 1 sekundę.

\end{description}


\paragraph{Bezpieczeństwo}
\label{\detokenize{rozdzial_2/rozdzial_22/index:bezpieczenstwo}}\begin{description}
\sphinxlineitem{\sphinxcode{\sphinxupquote{password\_encryption}}}
\sphinxAtStartPar
Algorytm szyfrowania haseł użytkowników.

\sphinxAtStartPar
Zalecana wartość:

\begin{sphinxVerbatim}[commandchars=\\\{\}]
\PYG{n+na}{password\PYGZus{}encryption}\PYG{+w}{ }\PYG{o}{=}\PYG{+w}{ }\PYG{l+s}{scram\PYGZhy{}sha\PYGZhy{}256}
\end{sphinxVerbatim}

\sphinxlineitem{\sphinxcode{\sphinxupquote{ssl}}}
\sphinxAtStartPar
Włącza obsługę połączeń SSL/TLS.

\sphinxAtStartPar
Dostępne wartości:
\begin{itemize}
\item {} 
\sphinxAtStartPar
\sphinxcode{\sphinxupquote{on}},

\item {} 
\sphinxAtStartPar
\sphinxcode{\sphinxupquote{off}}.

\end{itemize}

\sphinxlineitem{\sphinxcode{\sphinxupquote{ssl\_cert\_file}} / \sphinxcode{\sphinxupquote{ssl\_key\_file}}}
\sphinxAtStartPar
Ścieżki do certyfikatu i klucza prywatnego SSL.

\end{description}


\paragraph{Replikacja i WAL}
\label{\detokenize{rozdzial_2/rozdzial_22/index:replikacja-i-wal}}\begin{description}
\sphinxlineitem{\sphinxcode{\sphinxupquote{wal\_level}}}
\sphinxAtStartPar
Określa poziom szczegółowości Write\sphinxhyphen{}Ahead Log (WAL).

\sphinxAtStartPar
Dostępne wartości:
\begin{itemize}
\item {} 
\sphinxAtStartPar
\sphinxcode{\sphinxupquote{minimal}},

\item {} 
\sphinxAtStartPar
\sphinxcode{\sphinxupquote{replica}},

\item {} 
\sphinxAtStartPar
\sphinxcode{\sphinxupquote{logical}}.

\end{itemize}

\sphinxlineitem{\sphinxcode{\sphinxupquote{max\_wal\_size}}}
\sphinxAtStartPar
Maksymalny rozmiar plików WAL przed wykonaniem checkpointu.

\sphinxlineitem{\sphinxcode{\sphinxupquote{archive\_mode}}}
\sphinxAtStartPar
Włącza archiwizację WAL.

\sphinxlineitem{\sphinxcode{\sphinxupquote{archive\_command}}}
\sphinxAtStartPar
Polecenie wykonywane podczas archiwizacji plików WAL.

\end{description}


\subsubsection{Przykładowa konfiguracja}
\label{\detokenize{rozdzial_2/rozdzial_22/index:przykladowa-konfiguracja}}
\begin{sphinxVerbatim}[commandchars=\\\{\}]
\PYG{n+na}{listen\PYGZus{}addresses}\PYG{+w}{ }\PYG{o}{=}\PYG{+w}{ }\PYG{l+s}{\PYGZsq{}}\PYG{l+s}{*}\PYG{l+s}{\PYGZsq{}}
\PYG{n+na}{port}\PYG{+w}{ }\PYG{o}{=}\PYG{+w}{ }\PYG{l+s}{5432}
\PYG{n+na}{max\PYGZus{}connections}\PYG{+w}{ }\PYG{o}{=}\PYG{+w}{ }\PYG{l+s}{200}

\PYG{n+na}{shared\PYGZus{}buffers}\PYG{+w}{ }\PYG{o}{=}\PYG{+w}{ }\PYG{l+s}{1GB}
\PYG{n+na}{work\PYGZus{}mem}\PYG{+w}{ }\PYG{o}{=}\PYG{+w}{ }\PYG{l+s}{16MB}
\PYG{n+na}{maintenance\PYGZus{}work\PYGZus{}mem}\PYG{+w}{ }\PYG{o}{=}\PYG{+w}{ }\PYG{l+s}{256MB}

\PYG{n+na}{logging\PYGZus{}collector}\PYG{+w}{ }\PYG{o}{=}\PYG{+w}{ }\PYG{l+s}{on}
\PYG{n+na}{log\PYGZus{}directory}\PYG{+w}{ }\PYG{o}{=}\PYG{+w}{ }\PYG{l+s}{\PYGZsq{}}\PYG{l+s}{log}\PYG{l+s}{\PYGZsq{}}

\PYG{n+na}{ssl}\PYG{+w}{ }\PYG{o}{=}\PYG{+w}{ }\PYG{l+s}{on}
\PYG{n+na}{password\PYGZus{}encryption}\PYG{+w}{ }\PYG{o}{=}\PYG{+w}{ }\PYG{l+s}{scram\PYGZhy{}sha\PYGZhy{}256}

\PYG{n+na}{wal\PYGZus{}level}\PYG{+w}{ }\PYG{o}{=}\PYG{+w}{ }\PYG{l+s}{replica}
\PYG{n+na}{max\PYGZus{}wal\PYGZus{}size}\PYG{+w}{ }\PYG{o}{=}\PYG{+w}{ }\PYG{l+s}{2GB}
\end{sphinxVerbatim}


\subsubsection{Weryfikacja konfiguracji}
\label{\detokenize{rozdzial_2/rozdzial_22/index:weryfikacja-konfiguracji}}
\sphinxAtStartPar
Po zmianie parametrów konfigurację można przeładować bez restartu serwera:

\begin{sphinxVerbatim}[commandchars=\\\{\}]
SELECT\PYG{+w}{ }pg\PYGZus{}reload\PYGZus{}conf\PYG{o}{(}\PYG{o}{)}\PYG{p}{;}
\end{sphinxVerbatim}

\sphinxAtStartPar
Niektóre parametry wymagają pełnego restartu PostgreSQL.

\sphinxAtStartPar
Aktualne wartości parametrów można sprawdzić poleceniem:

\begin{sphinxVerbatim}[commandchars=\\\{\}]
\PYG{k}{SHOW}\PYG{+w}{ }\PYG{n}{shared\PYGZus{}buffers}\PYG{p}{;}
\PYG{k}{SHOW}\PYG{+w}{ }\PYG{n}{wal\PYGZus{}level}\PYG{p}{;}
\PYG{k}{SHOW}\PYG{+w}{ }\PYG{n}{max\PYGZus{}connections}\PYG{p}{;}
\end{sphinxVerbatim}


\subsection{pg\_hba.conf}
\label{\detokenize{rozdzial_2/rozdzial_22/index:pg-hba-conf}}
\sphinxAtStartPar
Plik \sphinxcode{\sphinxupquote{pg\_hba.conf}} (Host\sphinxhyphen{}Based Authentication) odpowiada za kontrolę dostępu do serwera PostgreSQL.

\sphinxAtStartPar
Definiuje:
\begin{itemize}
\item {} 
\sphinxAtStartPar
kto może połączyć się z bazą danych,

\item {} 
\sphinxAtStartPar
z jakiego adresu IP,

\item {} 
\sphinxAtStartPar
do jakiej bazy danych,

\item {} 
\sphinxAtStartPar
przy użyciu jakiej metody uwierzytelniania.

\end{itemize}

\sphinxAtStartPar
Lokalizację aktywnego pliku można sprawdzić poleceniem:

\begin{sphinxVerbatim}[commandchars=\\\{\}]
\PYG{k}{SHOW}\PYG{+w}{ }\PYG{n}{hba\PYGZus{}file}\PYG{p}{;}
\end{sphinxVerbatim}


\subsubsection{Składnia wpisu}
\label{\detokenize{rozdzial_2/rozdzial_22/index:skladnia-wpisu}}
\sphinxAtStartPar
Każdy wpis w \sphinxcode{\sphinxupquote{pg\_hba.conf}} ma następującą strukturę:

\begin{sphinxVerbatim}[commandchars=\\\{\}]
\PYG{n}{TYPE}    \PYG{n}{DATABASE}    \PYG{n}{USER}    \PYG{n}{ADDRESS}    \PYG{n}{METHOD}
\end{sphinxVerbatim}

\sphinxAtStartPar
Znaczenie pól:
\begin{description}
\sphinxlineitem{\sphinxcode{\sphinxupquote{TYPE}}}
\sphinxAtStartPar
Typ połączenia.

\sphinxAtStartPar
Dostępne wartości:
\begin{itemize}
\item {} 
\sphinxAtStartPar
\sphinxcode{\sphinxupquote{local}} \textendash{} połączenia lokalne (Unix socket),

\item {} 
\sphinxAtStartPar
\sphinxcode{\sphinxupquote{host}} \textendash{} połączenia TCP/IP,

\item {} 
\sphinxAtStartPar
\sphinxcode{\sphinxupquote{hostssl}} \textendash{} połączenia TCP/IP z SSL,

\item {} 
\sphinxAtStartPar
\sphinxcode{\sphinxupquote{hostnossl}} \textendash{} połączenia TCP/IP bez SSL.

\end{itemize}

\sphinxlineitem{\sphinxcode{\sphinxupquote{DATABASE}}}
\sphinxAtStartPar
Nazwa bazy danych, do której użytkownik może się połączyć.

\sphinxAtStartPar
Możliwe wartości:
\begin{itemize}
\item {} 
\sphinxAtStartPar
konkretna nazwa bazy,

\item {} 
\sphinxAtStartPar
\sphinxcode{\sphinxupquote{all}} \textendash{} wszystkie bazy,

\item {} 
\sphinxAtStartPar
\sphinxcode{\sphinxupquote{sameuser}} \textendash{} baza o nazwie zgodnej z użytkownikiem.

\end{itemize}

\sphinxlineitem{\sphinxcode{\sphinxupquote{USER}}}
\sphinxAtStartPar
Użytkownik lub rola PostgreSQL.

\sphinxlineitem{\sphinxcode{\sphinxupquote{ADDRESS}}}
\sphinxAtStartPar
Adres IP lub zakres sieci w notacji CIDR.

\sphinxAtStartPar
Przykłady:
\begin{itemize}
\item {} 
\sphinxAtStartPar
\sphinxcode{\sphinxupquote{127.0.0.1/32}},

\item {} 
\sphinxAtStartPar
\sphinxcode{\sphinxupquote{192.168.1.0/24}},

\item {} 
\sphinxAtStartPar
\sphinxcode{\sphinxupquote{0.0.0.0/0}} \textendash{} wszystkie adresy.

\end{itemize}

\sphinxlineitem{\sphinxcode{\sphinxupquote{METHOD}}}
\sphinxAtStartPar
Metoda uwierzytelniania użytkownika.

\end{description}


\subsubsection{Najczęściej używane metody uwierzytelniania}
\label{\detokenize{rozdzial_2/rozdzial_22/index:najczesciej-uzywane-metody-uwierzytelniania}}\begin{description}
\sphinxlineitem{\sphinxcode{\sphinxupquote{trust}}}
\sphinxAtStartPar
Zezwala na połączenie bez uwierzytelniania.

\sphinxAtStartPar
Zalecane wyłącznie w środowiskach testowych.

\sphinxlineitem{\sphinxcode{\sphinxupquote{reject}}}
\sphinxAtStartPar
Odrzuca połączenie.

\sphinxlineitem{\sphinxcode{\sphinxupquote{md5}}}
\sphinxAtStartPar
Uwierzytelnianie przy użyciu hasła MD5.

\sphinxAtStartPar
Metoda starsza i mniej bezpieczna niż SCRAM.

\sphinxlineitem{\sphinxcode{\sphinxupquote{scram\sphinxhyphen{}sha\sphinxhyphen{}256}}}
\sphinxAtStartPar
Nowoczesna metoda uwierzytelniania oparta na SCRAM.

\sphinxAtStartPar
Zalecana dla nowych instalacji PostgreSQL.

\sphinxlineitem{\sphinxcode{\sphinxupquote{peer}}}
\sphinxAtStartPar
Uwierzytelnianie na podstawie użytkownika systemowego.

\sphinxAtStartPar
Najczęściej używane dla lokalnych połączeń Linux/Unix.

\sphinxlineitem{\sphinxcode{\sphinxupquote{ident}}}
\sphinxAtStartPar
Mapowanie użytkownika systemowego przy użyciu \sphinxcode{\sphinxupquote{pg\_ident.conf}}.

\end{description}


\subsubsection{Przykładowa konfiguracja}
\label{\detokenize{rozdzial_2/rozdzial_22/index:id1}}
\begin{sphinxVerbatim}[commandchars=\\\{\}]
\PYG{c+c1}{\PYGZsh{} Połączenia lokalne}
\PYG{n+na}{local   all             postgres                        peer}

\PYG{c+c1}{\PYGZsh{} Połączenia localhost IPv4}
\PYG{n+na}{host    all             all         127.0.0.1/32        scram\PYGZhy{}sha\PYGZhy{}256}

\PYG{c+c1}{\PYGZsh{} Połączenia localhost IPv6}
\PYG{n+na}{host    all             all}\PYG{+w}{         }\PYG{o}{:}\PYG{l+s}{:1/128             scram\PYGZhy{}sha\PYGZhy{}256}

\PYG{c+c1}{\PYGZsh{} Sieć lokalna}
\PYG{n+na}{host    aplikacja       app\PYGZus{}user    192.168.1.0/24      scram\PYGZhy{}sha\PYGZhy{}256}
\end{sphinxVerbatim}


\subsubsection{Najlepsze praktyki}
\label{\detokenize{rozdzial_2/rozdzial_22/index:id2}}\begin{enumerate}
\sphinxsetlistlabels{\arabic}{enumi}{enumii}{}{.}%
\item {} 
\sphinxAtStartPar
Używaj \sphinxcode{\sphinxupquote{scram\sphinxhyphen{}sha\sphinxhyphen{}256}} zamiast \sphinxcode{\sphinxupquote{md5}}.

\item {} 
\sphinxAtStartPar
Ograniczaj dostęp do konkretnych adresów IP.

\item {} 
\sphinxAtStartPar
Nie używaj \sphinxcode{\sphinxupquote{trust}} w środowiskach produkcyjnych.

\item {} 
\sphinxAtStartPar
Stosuj \sphinxcode{\sphinxupquote{hostssl}} dla połączeń zdalnych.

\item {} 
\sphinxAtStartPar
Po zmianie konfiguracji wykonuj reload konfiguracji:

\end{enumerate}

\begin{sphinxVerbatim}[commandchars=\\\{\}]
\PYG{k}{SELECT}\PYG{+w}{ }\PYG{n}{pg\PYGZus{}reload\PYGZus{}conf}\PYG{p}{(}\PYG{p}{)}\PYG{p}{;}
\end{sphinxVerbatim}


\subsection{pg\_ident.conf}
\label{\detokenize{rozdzial_2/rozdzial_22/index:pg-ident-conf}}
\sphinxAtStartPar
Plik \sphinxcode{\sphinxupquote{pg\_ident.conf}} umożliwia mapowanie użytkowników systemowych na role PostgreSQL.

\sphinxAtStartPar
Najczęściej używany jest razem z metodami uwierzytelniania:
\begin{itemize}
\item {} 
\sphinxAtStartPar
\sphinxcode{\sphinxupquote{peer}},

\item {} 
\sphinxAtStartPar
\sphinxcode{\sphinxupquote{ident}}.

\end{itemize}

\sphinxAtStartPar
Lokalizację aktywnego pliku można sprawdzić poleceniem:

\begin{sphinxVerbatim}[commandchars=\\\{\}]
\PYG{k}{SHOW}\PYG{+w}{ }\PYG{n}{ident\PYGZus{}file}\PYG{p}{;}
\end{sphinxVerbatim}


\subsubsection{Zastosowanie}
\label{\detokenize{rozdzial_2/rozdzial_22/index:zastosowanie}}
\sphinxAtStartPar
Mechanizm mapowania pozwala określić, który użytkownik systemowy może zostać powiązany z konkretną rolą PostgreSQL.

\sphinxAtStartPar
Jest to szczególnie przydatne:
\begin{itemize}
\item {} 
\sphinxAtStartPar
w środowiskach Linux/Unix,

\item {} 
\sphinxAtStartPar
przy automatyzacji administracyjnej,

\item {} 
\sphinxAtStartPar
dla lokalnych połączeń bez użycia haseł.

\end{itemize}


\subsubsection{Składnia wpisu}
\label{\detokenize{rozdzial_2/rozdzial_22/index:id3}}
\begin{sphinxVerbatim}[commandchars=\\\{\}]
\PYG{n}{MAPNAME}    \PYG{n}{SYSTEM}\PYG{o}{\PYGZhy{}}\PYG{n}{USERNAME}    \PYG{n}{PG}\PYG{o}{\PYGZhy{}}\PYG{n}{USERNAME}
\end{sphinxVerbatim}

\sphinxAtStartPar
Znaczenie pól:
\begin{description}
\sphinxlineitem{\sphinxcode{\sphinxupquote{MAPNAME}}}
\sphinxAtStartPar
Nazwa mapowania wykorzystywana w \sphinxcode{\sphinxupquote{pg\_hba.conf}}.

\sphinxlineitem{\sphinxcode{\sphinxupquote{SYSTEM\sphinxhyphen{}USERNAME}}}
\sphinxAtStartPar
Użytkownik systemu operacyjnego.

\sphinxlineitem{\sphinxcode{\sphinxupquote{PG\sphinxhyphen{}USERNAME}}}
\sphinxAtStartPar
Rola PostgreSQL, na którą użytkownik ma zostać zmapowany.

\end{description}


\subsubsection{Przykładowa konfiguracja}
\label{\detokenize{rozdzial_2/rozdzial_22/index:id4}}
\begin{sphinxVerbatim}[commandchars=\\\{\}]
\PYG{c+c1}{\PYGZsh{} MAPNAME    SYSTEM\PYGZhy{}USERNAME    PG\PYGZhy{}USERNAME}
\PYG{n+na}{admin\PYGZus{}map    postgres           postgres}
\PYG{n+na}{admin\PYGZus{}map    deploy             db\PYGZus{}admin}
\PYG{n+na}{app\PYGZus{}map      appuser            aplikacja}
\end{sphinxVerbatim}

\sphinxAtStartPar
Przykład użycia w \sphinxcode{\sphinxupquote{pg\_hba.conf}}:

\begin{sphinxVerbatim}[commandchars=\\\{\}]
\PYG{n+na}{local   all     all     peer map}\PYG{o}{=}\PYG{l+s}{admin\PYGZus{}map}
\end{sphinxVerbatim}

\sphinxAtStartPar
W powyższym przykładzie użytkownicy systemowi zdefiniowani w \sphinxcode{\sphinxupquote{admin\_map}} mogą logować się jako odpowiadające im role PostgreSQL.


\subsubsection{Najlepsze praktyki}
\label{\detokenize{rozdzial_2/rozdzial_22/index:id5}}\begin{enumerate}
\sphinxsetlistlabels{\arabic}{enumi}{enumii}{}{.}%
\item {} 
\sphinxAtStartPar
Ograniczaj mapowania wyłącznie do wymaganych użytkowników.

\item {} 
\sphinxAtStartPar
Nie mapuj użytkowników systemowych na role superuser bez uzasadnienia.

\item {} 
\sphinxAtStartPar
Stosuj osobne role administracyjne i aplikacyjne.

\item {} 
\sphinxAtStartPar
Dokumentuj wszystkie mapowania użytkowników.

\item {} 
\sphinxAtStartPar
Regularnie przeglądaj konfigurację dostępu.

\end{enumerate}

\begin{sphinxadmonition}{note}{Opracowanie}

\sphinxAtStartPar
\sphinxstylestrong{Autor:} Michał Kraus
\sphinxstylestrong{Przedmiot:} Bazy danych
\sphinxstylestrong{Data:} maj 2026
\end{sphinxadmonition}

\sphinxstepscope


\section{Kontrola i konserwacja}
\label{\detokenize{rozdzial_2/rozdzial_23/index:kontrola-i-konserwacja}}\label{\detokenize{rozdzial_2/rozdzial_23/index::doc}}\begin{quote}\begin{description}
\sphinxlineitem{Autorzy}\begin{enumerate}
\sphinxsetlistlabels{\arabic}{enumi}{enumii}{}{.}%
\item {} 
\sphinxAtStartPar
Paweł Łoćwin

\item {} 
\sphinxAtStartPar
Paweł Łosowski

\end{enumerate}

\end{description}\end{quote}


\subsection{Wprowadzenie i planowanie konserwacji}
\label{\detokenize{rozdzial_2/rozdzial_23/index:wprowadzenie-i-planowanie-konserwacji}}
\sphinxAtStartPar
Współczesne systemy relacyjnych baz danych to wysoce skomplikowane środowiska, które wymagają ciągłego i proaktywnego nadzoru. Aby zagwarantować wysoką dostępność i niezawodność, niezbędne jest rygorystyczne planowanie procesów konserwacyjnych, które opiera się na odpowiedzi na trzy kluczowe pytania:
\begin{itemize}
\item {} 
\sphinxAtStartPar
\sphinxstylestrong{Kiedy?} Operacje inwazyjne (np. pełna przebudowa tabel i indeksów) powinny być planowane w tzw. oknach serwisowych, czyli w godzinach najniższego obciążenia systemu (zwykle w nocy lub w weekendy). Lżejsze prace konserwacyjne mogą odbywać się w tle.

\item {} 
\sphinxAtStartPar
\sphinxstylestrong{W jakim zakresie?} Należy precyzyjnie określić, czy konserwacji wymaga cała instancja, pojedyncza baza danych, czy tylko wytypowane, najszybciej rosnące tabele. Skupienie się na najbardziej obciążonych obszarach oszczędza zasoby sprzętowe.

\item {} 
\sphinxAtStartPar
\sphinxstylestrong{Jak?} Konserwacja powinna być maksymalnie zautomatyzowana (np. z wykorzystaniem systemowego demona CRON lub rozszerzenia \sphinxtitleref{pg\_cron}). Ręczne interwencje administratora powinny być zarezerwowane wyłącznie dla sytuacji awaryjnych i niestandardowych problemów wydajnościowych.

\end{itemize}


\subsection{Zarządzanie stanem serwera i sesjami}
\label{\detokenize{rozdzial_2/rozdzial_23/index:zarzadzanie-stanem-serwera-i-sesjami}}
\sphinxAtStartPar
Podstawą kontroli nad środowiskiem bazodanowym jest zarządzanie cyklem życia samej usługi. Uruchamianie, zatrzymywanie i restartowanie serwera bazy danych realizuje się najczęściej za pośrednictwem menedżera usług systemu operacyjnego (np. polecenia \sphinxtitleref{systemctl start/stop/restart postgresql}) lub za pomocą dedykowanego narzędzia wiersza poleceń \sphinxtitleref{pg\_ctl}.

\sphinxAtStartPar
Podczas wdrażania krytycznych aktualizacji lub w sytuacjach awaryjnych, konieczne jest sprawne zarządzanie ruchem użytkowników:
\begin{itemize}
\item {} 
\sphinxAtStartPar
\sphinxstylestrong{Zapobieganie nowym połączeniom:} Administrator może tymczasowo zablokować możliwość logowania się do konkretnej bazy danych poprzez zmianę jej metadanych. Pozwala to na „odcięcie” aplikacji bez wyłączania całego serwera.

\item {} 
\sphinxAtStartPar
\sphinxstylestrong{Ograniczanie i rozłączanie użytkowników:} Istniejące, aktywne sesje, które np. zawiesiły się lub blokują ważne operacje, można siłowo przerwać, zwalniając tym samym zablokowane zasoby.

\end{itemize}

\begin{sphinxVerbatim}[commandchars=\\\{\}]
\PYG{c+c1}{\PYGZhy{}\PYGZhy{} Zapobieganie nawiązywaniu nowych połączeń do bazy danych}
\PYG{k}{ALTER}\PYG{+w}{ }\PYG{k}{DATABASE}\PYG{+w}{ }\PYG{n}{biblioteka\PYGZus{}db}\PYG{+w}{ }\PYG{n}{ALLOW\PYGZus{}CONNECTIONS}\PYG{+w}{ }\PYG{o}{=}\PYG{+w}{ }\PYG{k}{false}\PYG{p}{;}

\PYG{c+c1}{\PYGZhy{}\PYGZhy{} Siłowe rozłączenie wszystkich aktywnych użytkowników (poza aktualną sesją)}
\PYG{k}{SELECT}\PYG{+w}{ }\PYG{n}{pg\PYGZus{}terminate\PYGZus{}backend}\PYG{p}{(}\PYG{n}{pid}\PYG{p}{)}
\PYG{k}{FROM}\PYG{+w}{ }\PYG{n}{pg\PYGZus{}stat\PYGZus{}activity}
\PYG{k}{WHERE}\PYG{+w}{ }\PYG{n}{datname}\PYG{+w}{ }\PYG{o}{=}\PYG{+w}{ }\PYG{l+s+s1}{\PYGZsq{}biblioteka\PYGZus{}db\PYGZsq{}}\PYG{+w}{ }\PYG{k}{AND}\PYG{+w}{ }\PYG{n}{pid}\PYG{+w}{ }\PYG{o}{\PYGZlt{}}\PYG{o}{\PYGZgt{}}\PYG{+w}{ }\PYG{n}{pg\PYGZus{}backend\PYGZus{}pid}\PYG{p}{(}\PYG{p}{)}\PYG{p}{;}
\end{sphinxVerbatim}


\subsection{Proces Vacuum}
\label{\detokenize{rozdzial_2/rozdzial_23/index:proces-vacuum}}
\sphinxAtStartPar
Architektura PostgreSQL opiera się na mechanizmie MVCC (Multi\sphinxhyphen{}Version Concurrency Control). Oznacza to, że instrukcje aktualizacji (UPDATE) i usuwania (DELETE) nie niszczą fizycznie starych danych, lecz tworzą tzw. „martwe krotki” (dead tuples). Do zarządzania nimi służy proces czyszczenia:
\begin{itemize}
\item {} 
\sphinxAtStartPar
\sphinxstylestrong{VACUUM:} Standardowa operacja zwalniająca miejsce po martwych wierszach, czyniąca je dostępnymi dla nowych danych. Nie zmniejsza ona fizycznego rozmiaru pliku na dysku i nie blokuje możliwości jednoczesnego odczytu i zapisu tabeli.

\item {} 
\sphinxAtStartPar
\sphinxstylestrong{VACUUM FULL:} Inwazyjna i ciężka operacja, która fizycznie przebudowuje tabelę, faktycznie zmniejszając jej rozmiar i zwracając wolne miejsce do systemu operacyjnego. Wymaga ekskluzywnej blokady (lock), całkowicie odcinając aplikację od modyfikowanej tabeli na czas trwania procesu.

\item {} 
\sphinxAtStartPar
\sphinxstylestrong{AUTOVACUUM:} Zautomatyzowany proces działający w tle, który regularnie monitoruje poziom zmian w tabelach i samoczynnie wyzwala standardowy Vacuum, zapobiegając niekontrolowanemu puchnięciu bazy danych (table bloat).

\end{itemize}


\subsection{Zarządzanie transakcjami i schematy (SCHEMA)}
\label{\detokenize{rozdzial_2/rozdzial_23/index:zarzadzanie-transakcjami-i-schematy-schema}}
\sphinxAtStartPar
Kolejnym aspektem administracji jest dbanie o logiczną strukturę bazy oraz spójność operacji.

\sphinxAtStartPar
\sphinxstylestrong{SCHEMA (Schematy)} to wirtualne przestrzenie nazw (przypominające katalogi w systemie operacyjnym), służące do logicznego grupowania tabel, widoków i funkcji. Ich główne zastosowania to:
* Izolacja danych różnych mikrousług lub aplikacji wewnątrz jednej wspólnej bazy danych.
* Implementacja architektury wielodostępowej (multi\sphinxhyphen{}tenant), gdzie każdy klient aplikacji obsługiwany jest w dedykowanym, oddzielnym schemacie.
* Ułatwione i zbiorcze zarządzanie uprawnieniami dostępu.

\sphinxAtStartPar
\sphinxstylestrong{Zarządzanie transakcjami} to proces kontrolowania instrukcji (BEGIN, COMMIT, ROLLBACK) w taki sposób, aby zachować zgodność z regułami ACID. Administrator musi zwracać szczególną uwagę na zapobieganie tzw. długim transakcjom (long\sphinxhyphen{}running transactions). Otwarta i porzucona transakcja powstrzymuje proces Vacuum przed usuwaniem martwych krotek, co prowadzi do drastycznego spadku wydajności całej bazy.

\begin{sphinxVerbatim}[commandchars=\\\{\}]
\PYG{c+c1}{\PYGZhy{}\PYGZhy{} Przykład logiki transakcyjnej z wykorzystaniem konkretnego schematu}
\PYG{k}{BEGIN}\PYG{p}{;}
\PYG{k}{INSERT}\PYG{+w}{ }\PYG{k}{INTO}\PYG{+w}{ }\PYG{n}{finanse}\PYG{p}{.}\PYG{n}{historia}\PYG{+w}{ }\PYG{p}{(}\PYG{n}{kwota}\PYG{p}{,}\PYG{+w}{ }\PYG{n}{opis}\PYG{p}{)}\PYG{+w}{ }\PYG{k}{VALUES}\PYG{+w}{ }\PYG{p}{(}\PYG{l+m+mi}{200}\PYG{p}{,}\PYG{+w}{ }\PYG{l+s+s1}{\PYGZsq{}Opłata serwisowa\PYGZsq{}}\PYG{p}{)}\PYG{p}{;}
\PYG{k}{UPDATE}\PYG{+w}{ }\PYG{n}{finanse}\PYG{p}{.}\PYG{n}{konta}\PYG{+w}{ }\PYG{k}{SET}\PYG{+w}{ }\PYG{n}{saldo}\PYG{+w}{ }\PYG{o}{=}\PYG{+w}{ }\PYG{n}{saldo}\PYG{+w}{ }\PYG{o}{\PYGZhy{}}\PYG{+w}{ }\PYG{l+m+mi}{200}\PYG{+w}{ }\PYG{k}{WHERE}\PYG{+w}{ }\PYG{n}{id\PYGZus{}klienta}\PYG{+w}{ }\PYG{o}{=}\PYG{+w}{ }\PYG{l+m+mi}{5}\PYG{p}{;}
\PYG{k}{COMMIT}\PYG{p}{;}
\end{sphinxVerbatim}


\subsection{Zarządzanie indeksami}
\label{\detokenize{rozdzial_2/rozdzial_23/index:zarzadzanie-indeksami}}
\sphinxAtStartPar
Indeksy są fundamentalne dla wydajnego wyszukiwania danych, jednak ich utrzymanie kosztuje — spowalniają one modyfikacje danych (DML) i konsumują miejsce na dysku.

\sphinxAtStartPar
Prawidłowe zarządzanie indeksami polega na stałym monitorowaniu ich wykorzystania. Należy identyfikować i usuwać indeksy, które nie są wykorzystywane przez planistę zapytań (tzw. unused indexes). Dodatkowo, podobnie jak tabele, indeksy ulegają fragmentacji. Z tego powodu administrator powinien regularnie planować operację \sphinxstylestrong{REINDEX}, która od nowa buduje strukturę drzewa indeksu, przywracając mu optymalny rozmiar i maksymalną wydajność. W systemach o wysokiej dostępności stosuje się przebudowę w trybie \sphinxtitleref{CONCURRENTLY}, która nie przerywa pracy aplikacji.


\subsection{Podsumowanie}
\label{\detokenize{rozdzial_2/rozdzial_23/index:podsumowanie}}
\sphinxAtStartPar
Prawidłowa kontrola i konserwacja środowiska bazodanowego to wielowymiarowy proces, wymagający głębokiego zrozumienia zarówno architektury logicznej, jak i fizycznej serwera.

\sphinxAtStartPar
Stabilność i wydajność systemu zależy w równej mierze od prawidłowego zarządzania stanem usługi i połączeniami, dogłębnego planowania okien serwisowych, jak i proaktywnego zarządzania mechanizmem Vacuum i transakcjami. Wykorzystanie schematów do strukturyzacji danych oraz regularne dbanie o kondycję indeksów pozwala na zbudowanie środowiska, które będzie skalowalne, przewidywalne i odporne na awarie w trudnych, produkcyjnych warunkach.
\begin{quote}\begin{description}
\sphinxlineitem{Autorzy}\begin{enumerate}
\sphinxsetlistlabels{\arabic}{enumi}{enumii}{}{.}%
\item {} 
\sphinxAtStartPar
Paweł Łoćwin

\item {} 
\sphinxAtStartPar
Paweł Łosowski

\end{enumerate}

\end{description}\end{quote}

\sphinxstepscope


\section{Monitorowanie i diagnostyka}
\label{\detokenize{rozdzial_2/rozdzial_24/index:monitorowanie-i-diagnostyka}}\label{\detokenize{rozdzial_2/rozdzial_24/index::doc}}\begin{quote}\begin{description}
\sphinxlineitem{Autorzy}\begin{enumerate}
\sphinxsetlistlabels{\arabic}{enumi}{enumii}{}{.}%
\item {} 
\sphinxAtStartPar
Oskar Wrona

\item {} 
\sphinxAtStartPar
Kamil Lewandowski

\item {} 
\sphinxAtStartPar
Adam Tarkowski

\end{enumerate}

\end{description}\end{quote}


\subsection{Wstęp}
\label{\detokenize{rozdzial_2/rozdzial_24/index:wstep}}
\sphinxAtStartPar
Monitorowanie i diagnostyka baz danych obejmują zestaw działań, które pozwalają administratorowi lub zespołowi utrzymania sprawdzić, czy system działa poprawnie, bezpiecznie i wydajnie. W praktyce nie chodzi jedynie o obserwowanie pojedynczych parametrów serwera, ale o łączenie informacji pochodzących z wielu warstw: samej bazy danych, systemu operacyjnego, aplikacji oraz narzędzi zewnętrznych. Dzięki temu można szybciej wykrywać przeciążenia, błędy konfiguracji, problemy z zapytaniami SQL, nieprawidłowe uprawnienia oraz awarie sprzętowe lub sieciowe.

\sphinxAtStartPar
W systemach bazodanowych szczególnie ważne jest to, że wiele problemów nie jest od razu widocznych dla użytkownika końcowego. Przykładowo zapytania mogą stopniowo wykonywać się coraz wolniej, liczba aktywnych sesji może rosnąć, a logi mogą zawierać powtarzające się ostrzeżenia. Bez monitorowania takie sygnały są łatwe do przeoczenia. Diagnostyka pozwala natomiast przejść od samego wykrycia problemu do ustalenia jego przyczyny, np. przez analizę sesji, blokad, planów zapytań, dzienników błędów lub zużycia zasobów systemowych.


\subsection{Sesje i użytkownicy}
\label{\detokenize{rozdzial_2/rozdzial_24/index:sesje-i-uzytkownicy}}
\sphinxAtStartPar
Jednym z podstawowych obszarów monitorowania bazy danych jest obserwowanie aktywnych sesji i użytkowników. Sesja oznacza połączenie użytkownika lub aplikacji z serwerem bazy danych. W ramach sesji wykonywane są zapytania, transakcje oraz operacje administracyjne. Analiza sesji pozwala sprawdzić między innymi, kto jest aktualnie połączony z bazą, z jakiego hosta pochodzi połączenie, jakie zapytanie jest wykonywane oraz czy sesja jest aktywna, bezczynna albo zablokowana.

\sphinxAtStartPar
W wielu systemach zarządzania bazami danych dostępne są specjalne widoki lub narzędzia administracyjne służące do obserwowania bieżącej aktywności. Na przykład PostgreSQL udostępnia mechanizmy monitorowania aktywności bazy danych oraz widoki statystyczne, takie jak \sphinxcode{\sphinxupquote{pg\_stat\_activity}}, które pozwalają sprawdzać informacje o procesach serwera i wykonywanych zapytaniach %
\begin{footnote}[1]\sphinxAtStartFootnote
PostgreSQL Documentation, \sphinxstyleemphasis{Monitoring Database Activity} oraz \sphinxstyleemphasis{The Cumulative Statistics System}, \sphinxurl{https://www.postgresql.org/docs/current/monitoring.html}
%
\end{footnote}. W MySQL podobną rolę może pełnić Performance Schema, które gromadzi dane diagnostyczne o pracy serwera i może być wykorzystywane do analizy wydajności %
\begin{footnote}[2]\sphinxAtStartFootnote
MySQL Documentation, \sphinxstyleemphasis{Using the Performance Schema to Diagnose Problems}, \sphinxurl{https://dev.mysql.com/doc/refman/8.2/en/performance-schema-examples.html}
%
\end{footnote}.

\sphinxAtStartPar
Monitorowanie sesji jest istotne zarówno z punktu widzenia wydajności, jak i bezpieczeństwa. Zbyt duża liczba jednoczesnych połączeń może prowadzić do przeciążenia serwera, a długo działające transakcje mogą blokować inne operacje. Z kolei nietypowe połączenia, np. z nieznanych adresów lub wykonywane poza normalnymi godzinami pracy, mogą wskazywać na błąd konfiguracji albo próbę nieautoryzowanego dostępu.


\subsection{Śledzenie dostępu użytkowników do tabel}
\label{\detokenize{rozdzial_2/rozdzial_24/index:sledzenie-dostepu-uzytkownikow-do-tabel}}
\sphinxAtStartPar
Kolejnym ważnym elementem jest kontrola tego, którzy użytkownicy uzyskują dostęp do określonych tabel i jakie operacje wykonują. W bazach danych przechowywane są często dane wrażliwe, dlatego sama informacja o tym, że użytkownik posiada konto w systemie, nie jest wystarczająca. Należy również wiedzieć, czy jego działania są zgodne z nadanymi uprawnieniami oraz z zasadami bezpieczeństwa obowiązującymi w organizacji.

\sphinxAtStartPar
Śledzenie dostępu może obejmować operacje odczytu danych, modyfikacji rekordów, usuwania danych, tworzenia nowych obiektów oraz zmian w uprawnieniach. W praktyce realizuje się to za pomocą mechanizmów audytu, dzienników zapytań, wyzwalaczy, narzędzi klasy Database Activity Monitoring lub funkcji wbudowanych w konkretny system bazodanowy. Przykładowo Oracle Database udostępnia mechanizmy audytu, a nowsze wersje zalecają korzystanie z unified auditing zamiast starszego audytu tradycyjnego %
\begin{footnote}[3]\sphinxAtStartFootnote
Oracle Documentation, \sphinxstyleemphasis{AUDIT\_TRAIL} oraz informacje o unified auditing, \sphinxurl{https://docs.oracle.com/en/database/oracle/oracle-database/21/refrn/AUDIT\_TRAIL.html}
%
\end{footnote}.

\sphinxAtStartPar
Takie śledzenie jest potrzebne nie tylko przy wykrywaniu incydentów bezpieczeństwa. Pomaga również w spełnianiu wymagań formalnych, analizie błędów użytkowników oraz odtwarzaniu przebiegu zdarzeń po awarii. Przykładowo, jeśli w tabeli pojawiły się nieprawidłowe dane, historia dostępu może pomóc ustalić, kiedy doszło do zmiany i która aplikacja lub który użytkownik ją wykonał.


\subsection{Błędy dziennika, logi i raporty}
\label{\detokenize{rozdzial_2/rozdzial_24/index:bledy-dziennika-logi-i-raporty}}
\sphinxAtStartPar
Dzienniki zdarzeń, czyli logi, są jednym z najważniejszych źródeł informacji diagnostycznych. Serwer bazy danych może zapisywać w nich komunikaty o błędach, ostrzeżeniach, nieudanych logowaniach, problemach z zapytaniami, restartach, zmianach konfiguracji lub przekroczeniu określonych progów wydajnościowych. Analiza logów pozwala wykrywać problemy, które nie zawsze są widoczne w aplikacji klienckiej.

\sphinxAtStartPar
W zależności od systemu bazodanowego logi mogą mieć różną postać. PostgreSQL obsługuje różne metody logowania komunikatów serwera, między innymi \sphinxcode{\sphinxupquote{stderr}}, \sphinxcode{\sphinxupquote{csvlog}}, \sphinxcode{\sphinxupquote{jsonlog}} oraz \sphinxcode{\sphinxupquote{syslog}}; w systemie Windows możliwe jest także użycie dziennika zdarzeń %
\begin{footnote}[4]\sphinxAtStartFootnote
PostgreSQL Documentation, \sphinxstyleemphasis{Error Reporting and Logging}, \sphinxurl{https://www.postgresql.org/docs/current/runtime-config-logging.html}
%
\end{footnote}. MySQL udostępnia między innymi general query log, który może rejestrować połączenia i zapytania otrzymywane przez serwer, oraz slow query log, używany do identyfikowania zapytań wykonujących się zbyt długo %
\begin{footnote}[5]\sphinxAtStartFootnote
MySQL Documentation, \sphinxstyleemphasis{The General Query Log} oraz \sphinxstyleemphasis{The Slow Query Log}, \sphinxurl{https://dev.mysql.com/doc/en/query-log.html}
%
\end{footnote}.

\sphinxAtStartPar
Raporty diagnostyczne są rozwinięciem surowych logów. Zamiast ręcznie analizować tysiące wpisów, administrator może korzystać z raportów podsumowujących liczbę błędów, najwolniejsze zapytania, częstotliwość nieudanych logowań lub zmiany obciążenia w czasie. Raporty są szczególnie przydatne w pracy zespołowej, ponieważ pozwalają dokumentować problemy i porównywać stan systemu przed oraz po wprowadzeniu zmian optymalizacyjnych.


\subsection{Monitorowanie na poziomie systemu operacyjnego}
\label{\detokenize{rozdzial_2/rozdzial_24/index:monitorowanie-na-poziomie-systemu-operacyjnego}}
\sphinxAtStartPar
Baza danych działa na konkretnym systemie operacyjnym i korzysta z jego zasobów. Dlatego diagnostyka nie może ograniczać się wyłącznie do samego silnika bazy danych. Wydajność zapytań może zależeć od procesora, pamięci RAM, operacji wejścia/wyjścia na dysku, sieci, liczby procesów oraz ogólnego obciążenia systemu.

\sphinxAtStartPar
W systemach Linux często wykorzystuje się narzędzia takie jak \sphinxcode{\sphinxupquote{iostat}}, \sphinxcode{\sphinxupquote{htop}} oraz \sphinxcode{\sphinxupquote{vmstat}}. Narzędzie \sphinxcode{\sphinxupquote{iostat}} pozwala obserwować wykorzystanie procesora i urządzeń wejścia/wyjścia, co jest szczególnie ważne przy bazach danych intensywnie korzystających z dysku. \sphinxcode{\sphinxupquote{htop}} umożliwia wygodne śledzenie procesów i zużycia zasobów w czasie rzeczywistym. \sphinxcode{\sphinxupquote{vmstat}} pokazuje między innymi informacje o pamięci, procesach, wymianie danych z dyskiem oraz obciążeniu procesora.

\sphinxAtStartPar
W systemach Microsoft Windows podobną rolę pełnią Menedżer zadań oraz Monitor wydajności. Pozwalają one obserwować zużycie procesora, pamięci, dysków, sieci oraz usług działających w tle. W środowiskach produkcyjnych takie dane są często zbierane automatycznie i zestawiane z metrykami pochodzącymi z bazy danych. Dzięki temu można odróżnić problem wynikający z nieoptymalnego zapytania SQL od problemu spowodowanego np. przeciążeniem dysku lub brakiem pamięci.


\subsection{Monitorowanie serwera}
\label{\detokenize{rozdzial_2/rozdzial_24/index:monitorowanie-serwera}}
\sphinxAtStartPar
Ostatnim poziomem jest monitorowanie całego serwera lub infrastruktury, w której działa baza danych. Obejmuje ono dostępność usług, czas odpowiedzi, zużycie zasobów, stan dysków, wykorzystanie sieci, działanie kopii zapasowych oraz wysyłanie powiadomień w przypadku awarii. W tym celu stosuje się narzędzia takie jak Nagios, Grafana, Prometheus, Zabbix lub inne systemy monitoringu.

\sphinxAtStartPar
Nagios jest przykładem narzędzia używanego do monitorowania usług i hostów, sprawdzania ich stanu oraz informowania administratorów o problemach %
\begin{footnote}[6]\sphinxAtStartFootnote
Nagios Open Source Documentation, \sphinxurl{https://www.nagios.org/documentation/}
%
\end{footnote}. Grafana z kolei służy przede wszystkim do wizualizacji danych monitorujących w formie dashboardów. Panele Grafany mogą prezentować metryki w postaci wykresów, tabel i innych wizualizacji, a moduł alertów pozwala reagować na określone zdarzenia lub przekroczenie progów %
\begin{footnote}[7]\sphinxAtStartFootnote
Grafana Documentation, \sphinxstyleemphasis{Dashboards overview} oraz \sphinxstyleemphasis{Grafana Alerting}, \sphinxurl{https://grafana.com/docs/grafana/latest/fundamentals/dashboards-overview/}
%
\end{footnote}.

\sphinxAtStartPar
Monitorowanie serwera jest szczególnie ważne w środowiskach, w których baza danych stanowi element większego systemu. Awaria bazy może wynikać nie tylko z błędu samego silnika, ale również z problemu sieciowego, zapełnionego dysku, braku pamięci, błędnej konfiguracji systemu lub nieudanego wdrożenia aplikacji. Centralne monitorowanie pozwala zebrać te informacje w jednym miejscu i szybciej wskazać źródło problemu.


\subsection{Znaczenie monitorowania i diagnostyki w bazach danych}
\label{\detokenize{rozdzial_2/rozdzial_24/index:znaczenie-monitorowania-i-diagnostyki-w-bazach-danych}}
\sphinxAtStartPar
Monitorowanie i diagnostyka są ważne, ponieważ wpływają na trzy kluczowe obszary: dostępność, wydajność i bezpieczeństwo danych. Dostępność oznacza, że baza danych jest osiągalna dla użytkowników i aplikacji. Wydajność oznacza, że zapytania wykonują się w akceptowalnym czasie. Bezpieczeństwo oznacza natomiast, że dostęp do danych jest kontrolowany, a podejrzane działania mogą zostać wykryte i przeanalizowane.

\sphinxAtStartPar
W dobrze utrzymanym środowisku bazodanowym monitorowanie powinno działać stale, a nie dopiero wtedy, gdy pojawi się awaria. Regularne obserwowanie sesji, logów, dostępu do tabel, parametrów systemu operacyjnego i stanu serwera pozwala wykrywać problemy na wczesnym etapie. Diagnostyka pozwala natomiast ustalić przyczyny tych problemów i zaplanować odpowiednie działania, takie jak optymalizacja zapytań, zmiana konfiguracji, zwiększenie zasobów, poprawa uprawnień lub wdrożenie dodatkowych zabezpieczeń.


\subsection{Podsumowanie}
\label{\detokenize{rozdzial_2/rozdzial_24/index:podsumowanie}}
\sphinxAtStartPar
Monitorowanie i diagnostyka baz danych tworzą wielowarstwowy proces kontroli działania systemu. Na poziomie bazy danych analizuje się sesje, użytkowników, zapytania, blokady, dostęp do tabel oraz logi. Na poziomie systemu operacyjnego obserwuje się zużycie procesora, pamięci, dysku i sieci. Na poziomie serwera wykorzystuje się narzędzia monitorujące, które zbierają metryki, tworzą wykresy oraz wysyłają alerty. Dopiero połączenie tych perspektyw daje pełny obraz stanu systemu i pozwala skutecznie reagować na błędy, spadki wydajności oraz potencjalne incydenty bezpieczeństwa.


\subsection{Źródła}
\label{\detokenize{rozdzial_2/rozdzial_24/index:zrodla}}
\sphinxstepscope


\section{Wydajność, skalowanie i replikacja}
\label{\detokenize{rozdzial_2/rozdzial_25/index:wydajnosc-skalowanie-i-replikacja}}\label{\detokenize{rozdzial_2/rozdzial_25/index::doc}}\begin{quote}\begin{description}
\sphinxlineitem{Autorzy}\begin{enumerate}
\sphinxsetlistlabels{\arabic}{enumi}{enumii}{}{.}%
\item {} 
\sphinxAtStartPar
Olaf Chomicki

\item {} 
\sphinxAtStartPar
Konrad Machowski

\item {} 
\sphinxAtStartPar
Wiktor Wydrzyński

\end{enumerate}

\end{description}\end{quote}


\subsection{Wydajność, skalowanie i replikacja}
\label{\detokenize{rozdzial_2/rozdzial_25/index:id1}}
\sphinxAtStartPar
Współczesne systemy bazodanowe muszą sprostać wymaganiom związanym z wykładniczym przyrostem wolumenu danych oraz stale rosnącą liczbą jednoczesnych użytkowników. W tym kontekście kluczowe staje się zrozumienie trzech fundamentalnych pojęć: wydajności, skalowania i replikacji. Wydajność odnosi się do szybkości przetwarzania pojedynczych transakcji i zapytań przy optymalnym wykorzystaniu dostępnych zasobów.

\sphinxAtStartPar
Gdy optymalizacja kodu przestaje wystarczać, konieczne staje się skalowanie \textendash{} pionowe (zwiększanie mocy pojedynczego serwera) lub poziome (rozpraszanie obciążenia na wiele maszyn). Replikacja z kolei stanowi most łączący skalowanie odczytów z zapewnieniem wysokiej dostępności i odporności systemu na awarie sprzętowe.


\subsection{Kontrola i buforowanie połączeń z bazą danych}
\label{\detokenize{rozdzial_2/rozdzial_25/index:kontrola-i-buforowanie-polaczen-z-baza-danych}}
\sphinxAtStartPar
Nawiązywanie nowego połączenia z bazą danych jest operacją niezwykle kosztowną procesorowo i czasowo, ponieważ wymaga uwierzytelnienia użytkownika, alokacji pamięci oraz ustanowienia sesji sieciowej. W systemach o dużym natężeniu ruchu brak kontroli nad liczbą połączeń prowadzi do szybkiego wyczerpania zasobów serwera.

\sphinxAtStartPar
Rozwiązaniem tego problemu jest buforowanie połączeń (Connection Pooling). Mechanizm ten polega na utrzymywaniu stałej puli otwartych, gotowych do użycia połączeń, które są wielokrotnie współdzielone przez aplikację. Narzędzia takie jak PgBouncer (dla PostgreSQL) czy wbudowane pule w serwerach aplikacji drastycznie zmniejszają narzut overhead, stabilizując czas reakcji bazy danych pod dużym obciążeniem.

\sphinxAtStartPar
Przykład konfiguracji PgBouncer (fragment pliku \sphinxcode{\sphinxupquote{pgbouncer.ini}}):

\begin{sphinxVerbatim}[commandchars=\\\{\}]
\PYG{k}{[databases]}
\PYG{n+na}{moja\PYGZus{}baza}\PYG{+w}{ }\PYG{o}{=}\PYG{+w}{ }\PYG{l+s}{host=127.0.0.1 port=5432 dbname=produkcja}

\PYG{k}{[pgbouncer]}
\PYG{n+na}{listen\PYGZus{}port}\PYG{+w}{ }\PYG{o}{=}\PYG{+w}{ }\PYG{l+s}{6432}
\PYG{n+na}{listen\PYGZus{}addr}\PYG{+w}{ }\PYG{o}{=}\PYG{+w}{ }\PYG{l+s}{*}
\PYG{n+na}{auth\PYGZus{}type}\PYG{+w}{ }\PYG{o}{=}\PYG{+w}{ }\PYG{l+s}{md5}
\PYG{n+na}{auth\PYGZus{}file}\PYG{+w}{ }\PYG{o}{=}\PYG{+w}{ }\PYG{l+s}{/etc/pgbouncer/userlist.txt}
\PYG{n+na}{pool\PYGZus{}mode}\PYG{+w}{ }\PYG{o}{=}\PYG{+w}{ }\PYG{l+s}{transaction}
\PYG{n+na}{max\PYGZus{}client\PYGZus{}conn}\PYG{+w}{ }\PYG{o}{=}\PYG{+w}{ }\PYG{l+s}{1000}
\PYG{n+na}{default\PYGZus{}pool\PYGZus{}size}\PYG{+w}{ }\PYG{o}{=}\PYG{+w}{ }\PYG{l+s}{20}
\end{sphinxVerbatim}


\subsection{INDEX i CLUSTER}
\label{\detokenize{rozdzial_2/rozdzial_25/index:index-i-cluster}}
\sphinxAtStartPar
Indeksowanie to podstawowa technika optymalizacji zapytań o charakterze odczytowym. Instrukcja \sphinxcode{\sphinxupquote{INDEX}} tworzy pomocniczą strukturę danych (najczęściej w formie B\sphinxhyphen{}drzewa), która pozwala silnikowi bazy danych na błyskawiczne odnalezienie żądanych wierszy bez konieczności kosztownego przeszukiwania całej tabeli (Full Table Scan).

\sphinxAtStartPar
Z kolei polecenie \sphinxcode{\sphinxupquote{CLUSTER}} idzie o krok dalej \textendash{} fizycznie reorganizuje strukturę danych na dysku w taki sposób, aby kolejność wierszy w tabeli dokładnie odpowiadała kolejności indeksu. Powoduje to, że powiązane rekordy są składowane obok siebie w tych samych blokach pamięci, co radykalnie przyspiesza zapytania zakresowe (Range Queries).

\sphinxAtStartPar
Przykład użycia poleceń w PostgreSQL:

\begin{sphinxVerbatim}[commandchars=\\\{\}]
\PYG{c+c1}{\PYGZhy{}\PYGZhy{} 1. Tworzenie standardowego indeksu B\PYGZhy{}drzewa na kolumnie data\PYGZus{}zamowienia}
\PYG{k}{CREATE}\PYG{+w}{ }\PYG{k}{INDEX}\PYG{+w}{ }\PYG{n}{idx\PYGZus{}zamowienia\PYGZus{}data}\PYG{+w}{ }\PYG{k}{ON}\PYG{+w}{ }\PYG{n}{zamowienia}\PYG{+w}{ }\PYG{p}{(}\PYG{n}{data\PYGZus{}zamowienia}\PYG{p}{)}\PYG{p}{;}

\PYG{c+c1}{\PYGZhy{}\PYGZhy{} 2. Fizyczne przemieszczenie danych na dysku według kolejności indeksu}
\PYG{k}{CLUSTER}\PYG{+w}{ }\PYG{n}{zamowienia}\PYG{+w}{ }\PYG{k}{USING}\PYG{+w}{ }\PYG{n}{idx\PYGZus{}zamowienia\PYGZus{}data}\PYG{p}{;}
\end{sphinxVerbatim}


\subsection{Rola i zastosowanie replikacji}
\label{\detokenize{rozdzial_2/rozdzial_25/index:rola-i-zastosowanie-replikacji}}
\sphinxAtStartPar
Replikacja polega na permanentnym kopiowaniu i synchronizowaniu danych z głównego węzła bazy danych (Primary/Master) na węzły pomocnicze (Replica/Slave). Jej rola w architekturze systemów IT jest wieloaspektowa. Po pierwsze, gwarantuje wysoką dostępność (High Availability) \textendash{} w przypadku fizycznej awarii Mastera, jedna z replik może automatycznie przejąć jego funkcje (proces Failover).

\sphinxAtStartPar
Po drugie, umożliwia skalowanie odczytów (Read Scalability). Przekierowanie operacji modyfikujących (INSERT, UPDATE) do węzła głównego, a operacji odczytu oraz generowania ciężkich raportów na repliki pozwala na znaczne odciążenie głównego serwera.


\subsection{Oprogramowanie i zaimplementowane mechanizmy replikacji}
\label{\detokenize{rozdzial_2/rozdzial_25/index:oprogramowanie-i-zaimplementowane-mechanizmy-replikacji}}
\sphinxAtStartPar
Mechanizmy replikacji mogą być realizowane na poziomie silnika bazy danych lub za pomocą zewnętrznego oprogramowania. Wyróżnia się dwa główne podejścia pod kątem transmisji danych: replikację opartą na logu (WAL \sphinxhyphen{} Write\sphinxhyphen{}Ahead Logging) oraz replikację logiczną. Ze względu na synchronizację, proces ten może przebiegać synchronicznie lub asynchronicznie.

\sphinxAtStartPar
Przykład konfiguracji replikacji logicznej (wbudowanej w PostgreSQL):

\begin{sphinxVerbatim}[commandchars=\\\{\}]
\PYG{c+c1}{\PYGZhy{}\PYGZhy{} KROK 1: Uruchomienie na węźle głównym (Primary) \PYGZhy{} zdefiniowanie publikacji}
\PYG{k}{CREATE}\PYG{+w}{ }\PYG{n}{PUBLICATION}\PYG{+w}{ }\PYG{n}{pub\PYGZus{}dane\PYGZus{}sprzedazy}\PYG{+w}{ }\PYG{k}{FOR}\PYG{+w}{ }\PYG{k}{TABLE}\PYG{+w}{ }\PYG{n}{zamowienia}\PYG{p}{,}\PYG{+w}{ }\PYG{n}{klienci}\PYG{p}{;}

\PYG{c+c1}{\PYGZhy{}\PYGZhy{} KROK 2: Uruchomienie na węźle pomocniczym (Replica) \PYGZhy{} subskrypcja danych}
\PYG{k}{CREATE}\PYG{+w}{ }\PYG{n}{SUBSCRIPTION}\PYG{+w}{ }\PYG{n}{sub\PYGZus{}dane\PYGZus{}sprzedazy}
\PYG{k}{CONNECTION}\PYG{+w}{ }\PYG{l+s+s1}{\PYGZsq{}host=192.168.1.50 port=5432 dbname=produkcja user=repl\PYGZus{}user password=secret\PYGZsq{}}
\PYG{n}{PUBLICATION}\PYG{+w}{ }\PYG{n}{pub\PYGZus{}dane\PYGZus{}sprzedazy}\PYG{p}{;}
\end{sphinxVerbatim}


\subsection{Limity systemu oraz ograniczanie dostępu użytkowników}
\label{\detokenize{rozdzial_2/rozdzial_25/index:limity-systemu-oraz-ograniczanie-dostepu-uzytkownikow}}
\sphinxAtStartPar
Bezpieczeństwo i stabilność bazy danych wymagają rygorystycznego zarządzania limitami systemowymi oraz uprawnieniami. Zbyt duża swoboda przyznana użytkownikom lub procesom aplikacyjnym może doprowadzić do celowego bądź przypadkowego unieruchomienia systemu.

\sphinxAtStartPar
Przykład nakładania restrykcji i limitów zasobów w PostgreSQL:

\begin{sphinxVerbatim}[commandchars=\\\{\}]
\PYG{c+c1}{\PYGZhy{}\PYGZhy{} 1. Ograniczenie maksymalnej liczby jednoczesnych połączeń dla danej roli}
\PYG{k}{ALTER}\PYG{+w}{ }\PYG{k}{ROLE}\PYG{+w}{ }\PYG{n}{rola\PYGZus{}analityk}\PYG{+w}{ }\PYG{k}{CONNECTION}\PYG{+w}{ }\PYG{k}{LIMIT}\PYG{+w}{ }\PYG{l+m+mi}{5}\PYG{p}{;}

\PYG{c+c1}{\PYGZhy{}\PYGZhy{} 2. Ustawienie maksymalnego czasu wykonywania pojedynczego zapytania (np. 30 sekund)}
\PYG{k}{ALTER}\PYG{+w}{ }\PYG{k}{ROLE}\PYG{+w}{ }\PYG{n}{rola\PYGZus{}analityk}\PYG{+w}{ }\PYG{k}{SET}\PYG{+w}{ }\PYG{n}{statement\PYGZus{}timeout}\PYG{+w}{ }\PYG{o}{=}\PYG{+w}{ }\PYG{l+s+s1}{\PYGZsq{}30s\PYGZsq{}}\PYG{p}{;}

\PYG{c+c1}{\PYGZhy{}\PYGZhy{} 3. Nadanie uprawnień wyłącznie do odczytu danych w wybranym schemacie}
\PYG{k}{REVOKE}\PYG{+w}{ }\PYG{k}{ALL}\PYG{+w}{ }\PYG{k}{ON}\PYG{+w}{ }\PYG{k}{ALL}\PYG{+w}{ }\PYG{n}{TABLES}\PYG{+w}{ }\PYG{k}{IN}\PYG{+w}{ }\PYG{k}{SCHEMA}\PYG{+w}{ }\PYG{k}{public}\PYG{+w}{ }\PYG{k}{FROM}\PYG{+w}{ }\PYG{n}{rola\PYGZus{}analityk}\PYG{p}{;}
\PYG{k}{GRANT}\PYG{+w}{ }\PYG{k}{SELECT}\PYG{+w}{ }\PYG{k}{ON}\PYG{+w}{ }\PYG{k}{ALL}\PYG{+w}{ }\PYG{n}{TABLES}\PYG{+w}{ }\PYG{k}{IN}\PYG{+w}{ }\PYG{k}{SCHEMA}\PYG{+w}{ }\PYG{k}{public}\PYG{+w}{ }\PYG{k}{TO}\PYG{+w}{ }\PYG{n}{rola\PYGZus{}analityk}\PYG{p}{;}
\end{sphinxVerbatim}


\subsection{Testy wydajności sprzętu na poziomie systemu operacyjnego}
\label{\detokenize{rozdzial_2/rozdzial_25/index:testy-wydajnosci-sprzetu-na-poziomie-systemu-operacyjnego}}
\sphinxAtStartPar
Przed wdrożeniem produkcyjnym bazy danych kluczowe jest przeprowadzenie niezależnych testów wydajnościowych podzespołów sprzętowych (Benchmarków) bezpośrednio na poziomie systemu operacyjnego, aby wykluczyć błędy konfiguracyjne środowiska.

\sphinxAtStartPar
Przykłady poleceń diagnostycznych i testowych w systemie Linux:

\begin{sphinxVerbatim}[commandchars=\\\{\}]
\PYG{c+c1}{\PYGZsh{} 1. Test wydajności procesora (CPU) za pomocą sysbench}
sysbench\PYG{+w}{ }cpu\PYG{+w}{ }\PYGZhy{}\PYGZhy{}cpu\PYGZhy{}max\PYGZhy{}prime\PYG{o}{=}\PYG{l+m}{20000}\PYG{+w}{ }run

\PYG{c+c1}{\PYGZsh{} 2. Test przepustowości i alokacji pamięci RAM}
sysbench\PYG{+w}{ }memory\PYG{+w}{ }\PYGZhy{}\PYGZhy{}memory\PYGZhy{}block\PYGZhy{}size\PYG{o}{=}1M\PYG{+w}{ }\PYGZhy{}\PYGZhy{}memory\PYGZhy{}total\PYGZhy{}size\PYG{o}{=}10G\PYG{+w}{ }run

\PYG{c+c1}{\PYGZsh{} 3. Zaawansowany test wydajności dysku (IOPS i opóźnienia) narzędziem fio}
\PYG{c+c1}{\PYGZsh{} Symulacja losowego odczytu/zapisu (Random R/W) specyficznego dla baz danych}
fio\PYG{+w}{ }\PYGZhy{}\PYGZhy{}name\PYG{o}{=}test\PYGZus{}bazy\PYG{+w}{ }\PYGZhy{}\PYGZhy{}ioengine\PYG{o}{=}libaio\PYG{+w}{ }\PYGZhy{}\PYGZhy{}rw\PYG{o}{=}randrw\PYG{+w}{ }\PYGZhy{}\PYGZhy{}bs\PYG{o}{=}4k\PYG{+w}{ }\PYG{l+s+se}{\PYGZbs{}}
\PYG{+w}{    }\PYGZhy{}\PYGZhy{}size\PYG{o}{=}2G\PYG{+w}{ }\PYGZhy{}\PYGZhy{}numjobs\PYG{o}{=}\PYG{l+m}{2}\PYG{+w}{ }\PYGZhy{}\PYGZhy{}runtime\PYG{o}{=}\PYG{l+m}{30}\PYG{+w}{ }\PYGZhy{}\PYGZhy{}time\PYGZus{}based\PYG{+w}{ }\PYG{l+s+se}{\PYGZbs{}}
\PYG{+w}{    }\PYGZhy{}\PYGZhy{}filename\PYG{o}{=}/var/lib/postgresql/test\PYGZus{}io.file\PYG{+w}{ }\PYGZhy{}\PYGZhy{}group\PYGZus{}reporting
\end{sphinxVerbatim}

\sphinxstepscope


\section{Partycjonowanie danych}
\label{\detokenize{rozdzial_2/rozdzial_26/index:partycjonowanie-danych}}\label{\detokenize{rozdzial_2/rozdzial_26/index::doc}}
\sphinxstepscope


\section{Bezpieczeństwo}
\label{\detokenize{rozdzial_2/rozdzial_27/index:bezpieczenstwo}}\label{\detokenize{rozdzial_2/rozdzial_27/index::doc}}\begin{quote}\begin{description}
\sphinxlineitem{Autor}
\sphinxAtStartPar
Mateusz

\end{description}\end{quote}


\subsection{Wstęp teoretyczny}
\label{\detokenize{rozdzial_2/rozdzial_27/index:wstep-teoretyczny}}
\sphinxAtStartPar
Bezpieczeństwo danych jest jednym z kluczowych aspektów zarządzania systemami baz danych. Obejmuje ono ochronę przed nieuprawnionym dostępem, złośliwą modyfikacją lub utratą informacji, a także zapewnienie poufności, integralności i ciągłej dostępności usług. W zależności od wybranej architektury systemu, mechanizmy bezpieczeństwa mogą być wbudowane bezpośrednio w silnik bazy danych lub w całości delegowane do systemu operacyjnego oraz aplikacji klienckiej.

\sphinxAtStartPar
W niniejszym podrozdziale dokonano szczegółowej analizy porównawczej mechanizmów bezpieczeństwa zaimplementowanych w sieciowym systemie \sphinxstylestrong{PostgreSQL} oraz wbudowanym silniku bazodanowym \sphinxstylestrong{SQLite}.


\subsection{Uwierzytelnianie i Kontrola Dostępu}
\label{\detokenize{rozdzial_2/rozdzial_27/index:uwierzytelnianie-i-kontrola-dostepu}}

\subsubsection{PostgreSQL}
\label{\detokenize{rozdzial_2/rozdzial_27/index:postgresql}}
\sphinxAtStartPar
PostgreSQL oferuje wysoce elastyczny podsystem uwierzytelniania, którego konfiguracja opiera się na pliku \sphinxcode{\sphinxupquote{pg\_hba.conf}} (Host\sphinxhyphen{}Based Authentication). Pozwala on na precyzyjne określenie, jaki użytkownik, z jakiego adresu IP i do jakiej bazy danych może się zalogować. System wspiera liczne zaawansowane metody weryfikacji tożsamości:
\begin{itemize}
\item {} 
\sphinxAtStartPar
\sphinxstylestrong{SCRAM\sphinxhyphen{}SHA\sphinxhyphen{}256}: Nowoczesny i bezpieczny domyślny mechanizm oparty na protokole wyzwanie\sphinxhyphen{}odpowiedź, uniemożliwiający podsłuchanie haseł.

\item {} 
\sphinxAtStartPar
\sphinxstylestrong{Uwierzytelnianie zewnętrzne}: Integracja z korporacyjnymi systemami zarządzania tożsamością, takimi jak LDAP, Kerberos, GSSAPI, PAM czy RADIUS.

\item {} 
\sphinxAtStartPar
\sphinxstylestrong{Certyfikaty SSL}: Możliwość uwierzytelniania klientów wyłącznie na podstawie ważnych certyfikatów kryptograficznych x509.

\end{itemize}


\subsubsection{SQLite}
\label{\detokenize{rozdzial_2/rozdzial_27/index:sqlite}}
\sphinxAtStartPar
W przeciwieństwie do systemów typu klient\sphinxhyphen{}serwer, SQLite nie posiada wbudowanego mechanizmu uwierzytelniania użytkowników. Wynika to bezpośrednio z jego architektury, gdzie baza danych jest po prostu pojedynczym plikiem na dysku.
\begin{itemize}
\item {} 
\sphinxAtStartPar
Kontrola dostępu do danych jest całkowicie zależna od uprawnień systemu operacyjnego.

\item {} 
\sphinxAtStartPar
Każdy proces lub użytkownik systemu operacyjnego, który posiada uprawnienia systemowe do odczytu i zapisu pliku bazy danych, ma automatycznie pełny dostęp do wszystkich danych i struktur w niej zawartych.

\item {} 
\sphinxAtStartPar
Istnieje oficjalne, opcjonalne rozszerzenie SQLite User Authentication, jednak nie jest ono częścią standardowej dystrybucji i zaprzecza podstawowej idei prostoty tego silnika.

\end{itemize}


\subsection{Autoryzacja i Uprawnienia}
\label{\detokenize{rozdzial_2/rozdzial_27/index:autoryzacja-i-uprawnienia}}

\subsubsection{PostgreSQL}
\label{\detokenize{rozdzial_2/rozdzial_27/index:id1}}
\sphinxAtStartPar
Autoryzacja w PostgreSQL realizowana jest poprzez model kontroli dostępu oparty na rolach (RBAC \textendash{} \sphinxstyleemphasis{Role\sphinxhyphen{}Based Access Control}). Za pomocą instrukcji \sphinxcode{\sphinxupquote{GRANT}} i \sphinxcode{\sphinxupquote{REVOKE}} administratorzy mogą szczegółowo definiować prawa:
\begin{itemize}
\item {} 
\sphinxAtStartPar
\sphinxstylestrong{Granularność uprawnień}: Uprawnienia mogą być nadawane na poziomie całych baz danych, schematów, tabel, widoków, a także konkretnych kolumn.

\item {} 
\sphinxAtStartPar
\sphinxstylestrong{Row\sphinxhyphen{}Level Security (RLS)}: Zaawansowany mechanizm polityk bezpieczeństwa, który filtruje zwracane wiersze w zależności od tożsamości zalogowanego użytkownika.

\end{itemize}


\subsubsection{SQLite}
\label{\detokenize{rozdzial_2/rozdzial_27/index:id2}}
\sphinxAtStartPar
Standardowy silnik SQLite nie obsługuje pojęcia ról ani wewnętrznych uprawnień.
\begin{itemize}
\item {} 
\sphinxAtStartPar
Po pomyślnym otwarciu pliku bazy danych, programista lub użytkownik może wykonać każdą operację (\sphinxcode{\sphinxupquote{SELECT}}, \sphinxcode{\sphinxupquote{INSERT}}, \sphinxcode{\sphinxupquote{DROP TABLE}}).

\item {} 
\sphinxAtStartPar
Jedyną formą ograniczenia uprawnień jest otwarcie bazy danych przez aplikację w trybie tylko do odczytu, użycie flagi wewnętrznej \sphinxcode{\sphinxupquote{SQLITE\_OPEN\_READONLY}} podczas inicjalizacji połączenia w kodzie źródłowym.

\end{itemize}


\subsection{Szyfrowanie Danych}
\label{\detokenize{rozdzial_2/rozdzial_27/index:szyfrowanie-danych}}

\subsubsection{PostgreSQL}
\label{\detokenize{rozdzial_2/rozdzial_27/index:id3}}
\sphinxAtStartPar
PostgreSQL zapewnia ochronę danych zarówno w trakcie ich przesyłania, jak i podczas przechowywania na nośnikach:
\begin{itemize}
\item {} 
\sphinxAtStartPar
\sphinxstylestrong{Szyfrowanie w locie (Data\sphinxhyphen{}in\sphinxhyphen{}Transit)}: Pełne i natywne wsparcie dla protokołu SSL/TLS, co zapobiega podsłuchiwaniu ruchu sieciowego między aplikacją a serwerem bazy danych.

\item {} 
\sphinxAtStartPar
\sphinxstylestrong{Szyfrowanie w spoczynku (Data\sphinxhyphen{}at\sphinxhyphen{}Rest)}: Standardowa wersja PostgreSQL nie posiada wbudowanego mechanizmu przezroczystego szyfrowania bazy danych (TDE). Bezpieczeństwo na dysku realizuje się za pomocą narzędzi systemowych lub poprzez wbudowane rozszerzenie \sphinxcode{\sphinxupquote{pgcrypto}}, pozwalające na szyfrowanie konkretnych wartości w kolumnach za pomocą algorytmu AES.

\end{itemize}


\subsubsection{SQLite}
\label{\detokenize{rozdzial_2/rozdzial_27/index:id4}}
\sphinxAtStartPar
Ponieważ cała baza danych SQLite mieści się w jednym pliku, fizyczne skopiowanie tego pliku oznacza kradzież wszystkich danych. Szyfrowanie staje się więc kluczowe:
\begin{itemize}
\item {} 
\sphinxAtStartPar
\sphinxstylestrong{Brak szyfrowania w standardzie}: Domyślnie pliki bazy danych są zapisywane jawnym tekstem.

\item {} 
\sphinxAtStartPar
\sphinxstylestrong{Szyfrowanie za pomocą rozszerzeń}: Aby zabezpieczyć bazę, stosuje się zewnętrzne biblioteki modyfikujące warstwę VFS (Virtual File System). Najpopularniejszym, otwartoźródłowym rozwiązaniem jest \sphinxstylestrong{SQLCipher} szyfrujący plik algorytmem AES\sphinxhyphen{}256 oraz oficjalne, komercyjne rozszerzenie \sphinxstylestrong{SEE} (\sphinxstyleemphasis{SQLite Encryption Extension}).

\end{itemize}


\subsection{Audyt i Logowanie}
\label{\detokenize{rozdzial_2/rozdzial_27/index:audyt-i-logowanie}}

\subsubsection{PostgreSQL}
\label{\detokenize{rozdzial_2/rozdzial_27/index:id5}}
\sphinxAtStartPar
System udostępnia zaawansowane mechanizmy śledzenia i rejestrowania operacji. Konfiguracja pozwala na logowanie błędów, prób uwierzytelnienia, a także pełnych zapytań SQL. W środowiskach wymagających rygorystycznego audytu stosuje się rozszerzenie \sphinxstylestrong{pgaudit}, które tworzy szczegółowe dzienniki zdarzeń zawierające informacje o tym, kto, kiedy i jakie dokładnie dane modyfikował lub przeglądał.


\subsubsection{SQLite}
\label{\detokenize{rozdzial_2/rozdzial_27/index:id6}}
\sphinxAtStartPar
SQLite nie generuje własnych plików dziennika zdarzeń przeznaczonych do audytu bezpieczeństwa. Śledzenie aktywności musi zostać zaimplementowane w kodzie aplikacji klienckiej przy użyciu funkcji API udostępnianych przez bibliotekę SQLite, takich jak mechanizmy przechwytujące zapytania: \sphinxcode{\sphinxupquote{sqlite3\_trace\_v2()}} lub filtry autoryzacyjne \sphinxcode{\sphinxupquote{sqlite3\_set\_authorizer()}}.


\subsection{Podsumowanie}
\label{\detokenize{rozdzial_2/rozdzial_27/index:podsumowanie}}
\sphinxAtStartPar
Podsumowując, podejście do bezpieczeństwa w systemach DBMS zależy w głównej mierze od ich założeń architektonicznych. PostgreSQL, jako zaawansowany system typu klient\sphinxhyphen{}serwer, oferuje kompleksowy i scentralizowany zestaw narzędzi bezpieczeństwa na poziomie silnika. Dzięki separacji procesów, zaawansowanej kontroli dostępu oraz natywnej obsłudze szyfrowania sieciowego, stanowi on optymalne rozwiązanie dla środowisk wielodostępnych i rozproszonych.

\sphinxAtStartPar
Z kolei SQLite jest minimalny i prosty. Jako bezserwerowy silnik osadzony, w całości deleguje zadania związane z uwierzytelnianiem i autoryzacją użytkowników do systemu operacyjnego lub samej aplikacji klienckiej. Bezpieczeństwo danych sprowadza się tu do ochrony fizycznego pliku bazy danych na dysku, a wdrożenie bardziej zaawansowanych mechanizmów, takich jak kryptograficzna ochrona danych w spoczynku, wymaga implementacji zewnętrznych nakładek programistycznych.

\sphinxstepscope


\section{Kopie zapasowe i odzyskiwanie danych}
\label{\detokenize{rozdzial_2/rozdzial_28/index:kopie-zapasowe-i-odzyskiwanie-danych}}\label{\detokenize{rozdzial_2/rozdzial_28/index::doc}}
\sphinxAtStartPar
\sphinxstylestrong{Autorzy:} Norbert Antonovitch, Michał Bednarczyk, Jan Balazs de Borbatviz


\subsection{Wstęp}
\label{\detokenize{rozdzial_2/rozdzial_28/index:wstep}}
\sphinxAtStartPar
Kopie zapasowe w PostgreSQL można podzielić na dwa główne nurty: kopie logiczne wykonywane narzędziami \sphinxcode{\sphinxupquote{pg\_dump}} i \sphinxcode{\sphinxupquote{pg\_dumpall}}  oraz kopie fizyczne całego klastra wykonywane m.in. przez \sphinxcode{\sphinxupquote{pg\_basebackup}} wraz z archiwizacją WAL dla odzyskiwania punktowego PITR (Point\sphinxhyphen{}in\sphinxhyphen{}Time Recovery). Kopie logiczne są wygodne do odtwarzania pojedynczych obiektów i pojedynczych baz, natomiast kopie fizyczne są podstawą pełnego odtworzenia instancji, tablespaces i odzyskiwania po awariach.


\subsection{Tworzenie kopii zapasowej całego systemu PostgreSQL \sphinxhyphen{} mechanizmy wbudowane}
\label{\detokenize{rozdzial_2/rozdzial_28/index:tworzenie-kopii-zapasowej-calego-systemu-postgresql-mechanizmy-wbudowane}}
\sphinxAtStartPar
Najprostszym wbudowanym sposobem wykonania kopii całego klastra jest \sphinxcode{\sphinxupquote{pg\_dumpall}}, które eksportuje wszystkie bazy w klastrze oraz obiekty globalne wspólne dla całej instancji, takie jak role i przestrzenie tabel. Taki zrzut ma postać tekstowego skryptu SQL i odtwarza się go zwykle przez \sphinxcode{\sphinxupquote{psql}}, ale metoda ta jest logiczna, więc nie daje możliwości odzyskiwania punktowego przez WAL.

\sphinxAtStartPar
Drugim, ważniejszym z punktu widzenia odporności na awarie mechanizmem jest \sphinxcode{\sphinxupquote{pg\_basebackup}}, które tworzy fizyczną kopię działającego klastra bez blokowania normalnej pracy klientów. Narzędzie to wykonuje kopię całego klastra, a nie pojedynczych baz czy tabel, może pracować w formacie katalogowym lub \sphinxcode{\sphinxupquote{tar}} i może dołączać wymagane pliki WAL metodą stream, fetch albo none. Dokumentacja podkreśla też, że taka kopia może służyć zarówno do PITR, jak i do zbudowania serwera standby.

\sphinxAtStartPar
Aby pełna kopia fizyczna była użyteczna w scenariuszu odtwarzania do wybranego momentu, należy włączyć archiwizację WAL przez ustawienie co najmniej \sphinxcode{\sphinxupquote{wal\_level = replica}}, \sphinxcode{\sphinxupquote{archive\_mode = on}} oraz poprawnego \sphinxcode{\sphinxupquote{archive\_command}} albo \sphinxcode{\sphinxupquote{archive\_library}}. PostgreSQL zapisuje każdą zmianę w logu WAL, a po odtworzeniu kopii bazowej można odtworzyć kolejne segmenty WAL, aby dojść do stanu bieżącego lub do wskazanego punktu w czasie.

\sphinxAtStartPar
Przykładowe polecenia:

\begin{sphinxVerbatim}[commandchars=\\\{\}]
\PYG{c+c1}{\PYGZsh{} logiczna kopia całego klastra}
pg\PYGZus{}dumpall\PYG{+w}{ }\PYGZhy{}U\PYG{+w}{ }postgres\PYG{+w}{ }\PYGZhy{}\PYGZhy{}clean\PYG{+w}{ }\PYGZhy{}\PYGZhy{}file\PYG{o}{=}cluster.sql

\PYG{c+c1}{\PYGZsh{} fizyczna kopia całego klastra z dołączeniem WAL}
pg\PYGZus{}basebackup\PYG{+w}{ }\PYGZhy{}h\PYG{+w}{ }dbserver\PYG{+w}{ }\PYGZhy{}D\PYG{+w}{ }/backup/base\PYG{+w}{ }\PYGZhy{}Fp\PYG{+w}{ }\PYGZhy{}X\PYG{+w}{ }stream\PYG{+w}{ }\PYGZhy{}P
\end{sphinxVerbatim}

\sphinxAtStartPar
W praktyce do dokumentacji laboratoryjnej warto zaznaczyć, że \sphinxcode{\sphinxupquote{pg\_dumpall}} lepiej nadaje się do migracji i prostych kopii administracyjnych, a \sphinxcode{\sphinxupquote{pg\_basebackup}} z archiwizacją WAL do scenariuszy disaster recovery.


\subsection{Tworzenie kopii zapasowych poszczególnych baz danych \sphinxhyphen{} mechanizmy wbudowane}
\label{\detokenize{rozdzial_2/rozdzial_28/index:tworzenie-kopii-zapasowych-poszczegolnych-baz-danych-mechanizmy-wbudowane}}
\sphinxAtStartPar
Do wykonania kopii pojedynczej bazy służy \sphinxcode{\sphinxupquote{pg\_dump}}, które tworzy spójny zrzut nawet wtedy, gdy baza jest równocześnie używana przez innych użytkowników, i nie blokuje zwykłych operacji odczytu ani zapisu. Narzędzie obsługuje format tekstowy, custom, directory oraz tar, przy czym formaty archiwalne współpracują z \sphinxcode{\sphinxupquote{pg\_restore}} i pozwalają na selektywne odtwarzanie obiektów.

\sphinxAtStartPar
Istotne jest to, że \sphinxcode{\sphinxupquote{pg\_dump}} wykonuje kopię tylko jednej bazy danych. Format custom (\sphinxcode{\sphinxupquote{\sphinxhyphen{}Fc}}) i directory (\sphinxcode{\sphinxupquote{\sphinxhyphen{}Fd}}) są najbardziej elastyczne, bo pozwalają wybierać odtwarzane elementy, zmieniać kolejność odtwarzania, a w części scenariuszy także korzystać z odtwarzania równoległego.

\sphinxAtStartPar
Przykładowe polecenia:

\begin{sphinxVerbatim}[commandchars=\\\{\}]
\PYG{c+c1}{\PYGZsh{} kopia pojedynczej bazy w formacie custom}
pg\PYGZus{}dump\PYG{+w}{ }\PYGZhy{}U\PYG{+w}{ }postgres\PYG{+w}{ }\PYGZhy{}d\PYG{+w}{ }labdb\PYG{+w}{ }\PYGZhy{}Fc\PYG{+w}{ }\PYGZhy{}f\PYG{+w}{ }/backup/labdb.dump

\PYG{c+c1}{\PYGZsh{} kopia pojedynczej bazy jako skrypt SQL}
pg\PYGZus{}dump\PYG{+w}{ }\PYGZhy{}U\PYG{+w}{ }postgres\PYG{+w}{ }\PYGZhy{}d\PYG{+w}{ }labdb\PYG{+w}{ }\PYGZhy{}Fp\PYG{+w}{ }\PYGZhy{}f\PYG{+w}{ }/backup/labdb.sql

\PYG{c+c1}{\PYGZsh{} odtworzenie archiwum custom}
pg\PYGZus{}restore\PYG{+w}{ }\PYGZhy{}d\PYG{+w}{ }labdb\PYGZus{}restored\PYG{+w}{ }/backup/labdb.dump
\end{sphinxVerbatim}

\sphinxAtStartPar
Jeżeli celem jest ochrona wybranych schematów lub tabel, \sphinxcode{\sphinxupquote{pg\_dump}} pozwala zawęzić zakres eksportu przez opcje takie jak \sphinxcode{\sphinxupquote{\sphinxhyphen{}n}} dla schematu i \sphinxcode{\sphinxupquote{\sphinxhyphen{}t}} dla tabeli. Trzeba jednak zaznaczyć, że wybiórczy zrzut nie gwarantuje kompletności zależności wszystkich obiektów w nowym, pustym środowisku, jeśli pominięto elementy, od których zależy odtwarzany obiekt.


\subsection{Odzyskiwanie usuniętych lub uszkodzonych danych}
\label{\detokenize{rozdzial_2/rozdzial_28/index:odzyskiwanie-usunietych-lub-uszkodzonych-danych}}
\sphinxAtStartPar
Sposób odzyskiwania zależy od tego, czy problem dotyczy pojedynczego obiektu logicznego, całej bazy, przestrzeni tabel, czy fizycznego uszkodzenia danych. PostgreSQL rozróżnia w praktyce odzyskiwanie logiczne z plików dump oraz odzyskiwanie fizyczne z kopii bazowej i archiwum WAL.


\subsubsection{Odtwarzanie tabel i innych obiektów logicznych}
\label{\detokenize{rozdzial_2/rozdzial_28/index:odtwarzanie-tabel-i-innych-obiektow-logicznych}}
\sphinxAtStartPar
Jeżeli wcześniej wykonano kopię \sphinxcode{\sphinxupquote{pg\_dump}} w formacie archiwalnym, można odtworzyć pojedynczą tabelę lub inny obiekt selektywnie przy pomocy \sphinxcode{\sphinxupquote{pg\_restore}}, bez przywracania całej bazy. To jest najprostszy wariant odzyskania przypadkowo usuniętej tabeli, o ile odpowiedni obiekt znajdował się w logicznej kopii zapasowej.

\sphinxAtStartPar
Przykład:

\begin{sphinxVerbatim}[commandchars=\\\{\}]
pg\PYGZus{}restore\PYG{+w}{ }\PYGZhy{}d\PYG{+w}{ }labdb\PYG{+w}{ }\PYGZhy{}t\PYG{+w}{ }public.wyniki\PYG{+w}{ }/backup/labdb.dump
\end{sphinxVerbatim}


\subsubsection{Odtwarzanie usuniętej bazy danych}
\label{\detokenize{rozdzial_2/rozdzial_28/index:odtwarzanie-usunietej-bazy-danych}}
\sphinxAtStartPar
Usuniętą bazę można najłatwiej odtworzyć z wcześniejszego zrzutu \sphinxcode{\sphinxupquote{pg\_dump}} tej konkretnej bazy albo z \sphinxcode{\sphinxupquote{pg\_dumpall}}, jeżeli wykonywano kopię całego klastra. W przypadku wymogu odtworzenia dokładnie do chwili sprzed usunięcia, potrzebna jest kopia fizyczna oraz archiwum WAL, ponieważ same zrzuty logiczne nie są częścią rozwiązania continuous archiving i nie wspierają PITR.


\subsubsection{Odtwarzanie przestrzeni tabel i pełnego klastra}
\label{\detokenize{rozdzial_2/rozdzial_28/index:odtwarzanie-przestrzeni-tabel-i-pelnego-klastra}}
\sphinxAtStartPar
Przestrzenie tabel są obiektami globalnymi klastra, dlatego ich definicje są uwzględniane przez \sphinxcode{\sphinxupquote{pg\_dumpall}}, a przy kopiach fizycznych \sphinxcode{\sphinxupquote{pg\_basebackup}} zachowuje także ich strukturę i mapowanie. Dokumentacja \sphinxcode{\sphinxupquote{pg\_basebackup}} wskazuje, że przy pracy z tablespaces można użyć mapowania \sphinxcode{\sphinxupquote{\sphinxhyphen{}\sphinxhyphen{}tablespace\sphinxhyphen{}mapping}}, a przy odtwarzaniu trzeba zweryfikować poprawność dowiązań symbolicznych i lokalizacji danych.

\sphinxAtStartPar
Przy fizycznym uszkodzeniu danych, katalogu PGDATA albo tabelpaces standardowa procedura obejmuje odtworzenie kopii bazowej, usunięcie starych plików danych, odtworzenie katalogów danych i tablespaces, skonfigurowanie \sphinxcode{\sphinxupquote{restore\_command}}, utworzenie pliku \sphinxcode{\sphinxupquote{recovery.signal}}, a następnie ponowne uruchomienie serwera w trybie recovery. PostgreSQL podczas odtwarzania odczytuje wymagane segmenty WAL i może dojść do stanu bieżącego albo zatrzymać się na wybranym punkcie odzyskiwania.


\subsubsection{PITR i odzyskiwanie stanu sprzed błędu}
\label{\detokenize{rozdzial_2/rozdzial_28/index:pitr-i-odzyskiwanie-stanu-sprzed-bledu}}
\sphinxAtStartPar
Najważniejszym mechanizmem odzyskiwania po usunięciu danych operacyjnych jest PITR, czyli przywrócenie systemu do chwili sprzed błędnej operacji, na przykład \sphinxcode{\sphinxupquote{DROP TABLE}} lub \sphinxcode{\sphinxupquote{DROP DATABASE}}. W takim scenariuszu odtwarza się kopię bazową, konfiguruje \sphinxcode{\sphinxupquote{restore\_command}}, a następnie ustawia cel odzyskiwania, np. przez czas, nazwany punkt przywracania lub identyfikator transakcji.

\sphinxAtStartPar
To rozwiązanie ma jednak ważne ograniczenie: nie służy do odzyskania pojedynczej tabeli „w miejscu” bez wpływu na resztę systemu, lecz do odtworzenia całego klastra do wcześniejszego stanu. Dlatego w praktyce laboratoryjnej często łączy się oba podejścia: regularne \sphinxcode{\sphinxupquote{pg\_dump}} dla obiektów logicznych i \sphinxcode{\sphinxupquote{pg\_basebackup}} z archiwizacją WAL dla pełnego recovery.


\subsection{Dedykowane oprogramowanie i skrypty zewnętrzne do automatyzacji}
\label{\detokenize{rozdzial_2/rozdzial_28/index:dedykowane-oprogramowanie-i-skrypty-zewnetrzne-do-automatyzacji}}
\sphinxAtStartPar
Wbudowane narzędzia PostgreSQL są wystarczające do ręcznego wykonywania kopii, ale w środowisku produkcyjnym często wykorzystuje się oprogramowanie automatyzujące harmonogramy, retencję, walidację i odtwarzanie. Dwa najczęściej wskazywane rozwiązania to Barman i pgBackRest, przy czym dostępna dokumentacja Barman bezpośrednio opisuje obsługę pełnych i przyrostowych kopii, archiwizacji WAL oraz automatycznego uruchamiania backupów.

\sphinxAtStartPar
Barman wykonuje kopie całego serwera PostgreSQL i wymaga poprawnej archiwizacji WAL, oferując różne metody kopii, w tym streaming przez \sphinxcode{\sphinxupquote{pg\_basebackup}}, \sphinxcode{\sphinxupquote{rsync}} oraz backupy snapshotowe w chmurze. Dokumentacja podaje też wsparcie dla backupów przyrostowych, limitowania pasma, kompresji i pracy z harmonogramem przez zewnętrzne mechanizmy systemowe, np. cron.

\sphinxAtStartPar
Przykładowe zastosowania narzędzi zewnętrznych:
\begin{itemize}
\item {} 
\sphinxAtStartPar
\sphinxstylestrong{Barman:} centralne repozytorium backupów, archiwizacja WAL, backupy pełne i przyrostowe, odzyskiwanie całych instancji.

\item {} 
\sphinxAtStartPar
\sphinxstylestrong{pgBackRest:} popularne narzędzie do wydajnych backupów i retencji, często używane w środowiskach HA oraz większych instalacjach PostgreSQL.

\item {} 
\sphinxAtStartPar
\sphinxstylestrong{Własne skrypty Bash/PowerShell:} wywołanie \sphinxcode{\sphinxupquote{pg\_dump}}, \sphinxcode{\sphinxupquote{pg\_dumpall}} lub \sphinxcode{\sphinxupquote{pg\_basebackup}}, kompresja plików, rotacja katalogów, logowanie i wysyłka wyników przez e\sphinxhyphen{}mail lub do systemu monitoringu.

\end{itemize}

\sphinxAtStartPar
Przykładowy prosty skrypt automatyzujący kopię jednej bazy:

\begin{sphinxVerbatim}[commandchars=\\\{\}]
\PYG{c+ch}{\PYGZsh{}!/bin/bash}
\PYG{n+nv}{DATA}\PYG{o}{=}\PYG{k}{\PYGZdl{}(}date\PYG{+w}{ }+\PYGZpc{}F\PYGZus{}\PYGZpc{}H\PYGZhy{}\PYGZpc{}M\PYG{k}{)}
\PYG{n+nv}{KATALOG}\PYG{o}{=}/backup/postgres
\PYG{n+nv}{BAZA}\PYG{o}{=}labdb

mkdir\PYG{+w}{ }\PYGZhy{}p\PYG{+w}{ }\PYG{l+s+s2}{\PYGZdq{}}\PYG{n+nv}{\PYGZdl{}KATALOG}\PYG{l+s+s2}{\PYGZdq{}}
pg\PYGZus{}dump\PYG{+w}{ }\PYGZhy{}U\PYG{+w}{ }postgres\PYG{+w}{ }\PYGZhy{}d\PYG{+w}{ }\PYG{l+s+s2}{\PYGZdq{}}\PYG{n+nv}{\PYGZdl{}BAZA}\PYG{l+s+s2}{\PYGZdq{}}\PYG{+w}{ }\PYGZhy{}Fc\PYG{+w}{ }\PYGZhy{}f\PYG{+w}{ }\PYG{l+s+s2}{\PYGZdq{}}\PYG{n+nv}{\PYGZdl{}KATALOG}\PYG{l+s+s2}{/}\PYG{l+s+si}{\PYGZdl{}\PYGZob{}}\PYG{n+nv}{BAZA}\PYG{l+s+si}{\PYGZcb{}}\PYG{l+s+s2}{\PYGZus{}}\PYG{l+s+si}{\PYGZdl{}\PYGZob{}}\PYG{n+nv}{DATA}\PYG{l+s+si}{\PYGZcb{}}\PYG{l+s+s2}{.dump}\PYG{l+s+s2}{\PYGZdq{}}
find\PYG{+w}{ }\PYG{l+s+s2}{\PYGZdq{}}\PYG{n+nv}{\PYGZdl{}KATALOG}\PYG{l+s+s2}{\PYGZdq{}}\PYG{+w}{ }\PYGZhy{}type\PYG{+w}{ }f\PYG{+w}{ }\PYGZhy{}name\PYG{+w}{ }\PYG{l+s+s2}{\PYGZdq{}}\PYG{l+s+si}{\PYGZdl{}\PYGZob{}}\PYG{n+nv}{BAZA}\PYG{l+s+si}{\PYGZcb{}}\PYG{l+s+s2}{\PYGZus{}*.dump}\PYG{l+s+s2}{\PYGZdq{}}\PYG{+w}{ }\PYGZhy{}mtime\PYG{+w}{ }+7\PYG{+w}{ }\PYGZhy{}delete
\end{sphinxVerbatim}

\sphinxAtStartPar
Taki skrypt jest prosty, ale nie zapewnia pełnego disaster recovery, bo nie obejmuje archiwizacji WAL i nie chroni całego klastra. Z tego powodu dla systemów krytycznych bardziej właściwe są narzędzia klasy Barman lub pgBackRest, które porządkują pełny proces backupu, retencji, testów i odtwarzania.


\subsection{Podsumowanie}
\label{\detokenize{rozdzial_2/rozdzial_28/index:podsumowanie}}
\sphinxAtStartPar
W rozdziale warto wyraźnie rozróżnić kopie logiczne i fizyczne, bo odpowiadają one na inne potrzeby eksploatacyjne. \sphinxcode{\sphinxupquote{pg\_dump}} i \sphinxcode{\sphinxupquote{pg\_dumpall}} są najlepsze do prostych eksportów i selektywnego odtwarzania, natomiast \sphinxcode{\sphinxupquote{pg\_basebackup}} z archiwizacją WAL jest podstawą odzyskiwania po poważnej awarii i realizacji PITR.

\sphinxstepscope


\chapter{Planowanie baz danych i tworzenie dokumentacji}
\label{\detokenize{rozdzial_3/index:planowanie-baz-danych-i-tworzenie-dokumentacji}}\label{\detokenize{rozdzial_3/index::doc}}
\sphinxAtStartPar
Autor: Mateusz Gałecki


\section{Wybór zagadnienia i opis procesów}
\label{\detokenize{rozdzial_3/index:wybor-zagadnienia-i-opis-procesow}}
\sphinxAtStartPar
Do stworzenia wybrałem baze danych obsługującą system kina.

\sphinxAtStartPar
\sphinxstylestrong{Opis głównych procesów:}
\begin{itemize}
\item {} 
\sphinxAtStartPar
\sphinxstylestrong{Zarządzanie katalogiem filmów:} Wprowadzanie danych filmowych

\item {} 
\sphinxAtStartPar
\sphinxstylestrong{Zarządzanie infrastrukturą:} Dane o salach kinowych ich pojemności i sprzedaży biletów

\item {} 
\sphinxAtStartPar
\sphinxstylestrong{Planowanie repertuaru:} Przypisanie filmów dla konkretnych sal

\end{itemize}


\section{Prototypowe dane}
\label{\detokenize{rozdzial_3/index:prototypowe-dane}}
\sphinxAtStartPar
Prototypowe pliki w formatach CSV i JSON.

\sphinxAtStartPar
\sphinxstylestrong{Format CSV:}

\begin{sphinxVerbatim}[commandchars=\\\{\}]
tytul\PYGZus{}filmu,rok\PYGZus{}produkcji,czas\PYGZus{}trwania,nazwa\PYGZus{}sali,pojemnosc,czas\PYGZus{}rozpoczecia
Matrix,1999,136,Sala 1,120,2026\PYGZhy{}06\PYGZhy{}25 18:00:00
\end{sphinxVerbatim}

\sphinxAtStartPar
\sphinxstylestrong{Format JSON:}

\begin{sphinxVerbatim}[commandchars=\\\{\}]
\PYG{p}{\PYGZob{}}
\PYG{+w}{  }\PYG{n+nt}{\PYGZdq{}tytul\PYGZdq{}}\PYG{p}{:}\PYG{+w}{ }\PYG{l+s+s2}{\PYGZdq{}Matrix\PYGZdq{}}\PYG{p}{,}
\PYG{+w}{  }\PYG{n+nt}{\PYGZdq{}rok\PYGZus{}produkcji\PYGZdq{}}\PYG{p}{:}\PYG{+w}{ }\PYG{l+m+mi}{1999}\PYG{p}{,}
\PYG{+w}{  }\PYG{n+nt}{\PYGZdq{}czas\PYGZus{}trwania\PYGZdq{}}\PYG{p}{:}\PYG{+w}{ }\PYG{l+m+mi}{136}\PYG{p}{,}
\PYG{+w}{  }\PYG{n+nt}{\PYGZdq{}nazwa\PYGZus{}sali\PYGZdq{}}\PYG{p}{:}\PYG{+w}{ }\PYG{l+s+s2}{\PYGZdq{}Sala 1\PYGZdq{}}\PYG{p}{,}
\PYG{+w}{  }\PYG{n+nt}{\PYGZdq{}pojemnosc\PYGZdq{}}\PYG{p}{:}\PYG{+w}{ }\PYG{l+m+mi}{120}\PYG{p}{,}
\PYG{+w}{  }\PYG{n+nt}{\PYGZdq{}czas\PYGZus{}rozpoczecia\PYGZdq{}}\PYG{p}{:}\PYG{+w}{ }\PYG{l+s+s2}{\PYGZdq{}2026\PYGZhy{}06\PYGZhy{}25 18:00:00\PYGZdq{}}
\PYG{p}{\PYGZcb{}}
\end{sphinxVerbatim}


\section{Model konceptualny}
\label{\detokenize{rozdzial_3/index:model-konceptualny}}
\sphinxAtStartPar
Podzieliłem to na następujące elementy:
\begin{itemize}
\item {} 
\sphinxAtStartPar
\sphinxstylestrong{Encje silne:} Film, Sala (ich istnienie jest niezależne od innych)

\item {} 
\sphinxAtStartPar
\sphinxstylestrong{Encje słabe:} Seans (jej istnienie jest uzależnione od innych encji)

\item {} 
\sphinxAtStartPar
\sphinxstylestrong{Związki niepoprawne:} pozbyłem się związków wiele\sphinxhyphen{}do\sphinxhyphen{}wielu pomiędzy filmem a salą.

\end{itemize}

\noindent{\hspace*{\fill}\sphinxincludegraphics{{diagram_chen}.png}\hspace*{\fill}}


\section{Model logiczny i normalizacja}
\label{\detokenize{rozdzial_3/index:model-logiczny-i-normalizacja}}
\sphinxAtStartPar
Dane z prototypu poddałem normalizacji osiągając 3 postać normalną (3NF).
\begin{enumerate}
\sphinxsetlistlabels{\arabic}{enumi}{enumii}{}{.}%
\item {} 
\sphinxAtStartPar
\sphinxstylestrong{1NF:} każda informacja znajduje się w osobnej kolumnie

\item {} 
\sphinxAtStartPar
\sphinxstylestrong{2NF i 3NF:} Usunąłem zależności częściowe i przechodnie.

\end{enumerate}

\sphinxAtStartPar
Związki miedzy znormalizowanymi tabelami maja postać 1:N
\begin{itemize}
\item {} 
\sphinxAtStartPar
Jeden FILM może posiadać wiele SEANSÓW.

\item {} 
\sphinxAtStartPar
Jedna SALA może posiadać wiele SEANSÓW.

\end{itemize}

\noindent{\hspace*{\fill}\sphinxincludegraphics{{diagram_erd}.png}\hspace*{\fill}}


\section{Model fizyczny bazy danych}
\label{\detokenize{rozdzial_3/index:model-fizyczny-bazy-danych}}
\sphinxAtStartPar
Model bazy danych z podziałem na PostgreSQL i SQLite

\sphinxAtStartPar
\sphinxstylestrong{PostgreSQL}


\begin{savenotes}\sphinxattablestart
\sphinxthistablewithglobalstyle
\centering
\sphinxcapstartof{table}
\sphinxthecaptionisattop
\sphinxcaption{Model fizyczny bazy danych (PostgreSQL)}\label{\detokenize{rozdzial_3/index:id1}}
\sphinxaftertopcaption
\begin{tabular}[t]{\X{20}{100}\X{25}{100}\X{25}{100}\X{30}{100}}
\sphinxtoprule
\begin{varwidth}[t]{\sphinxcolwidth{1}{4}}
\sphinxstyletheadfamily \sphinxAtStartPar
Tabela
\sphinxbeforeendvarwidth
\end{varwidth}%
&\begin{varwidth}[t]{\sphinxcolwidth{1}{4}}
\sphinxstyletheadfamily \sphinxAtStartPar
Kolumna
\sphinxbeforeendvarwidth
\end{varwidth}%
&\begin{varwidth}[t]{\sphinxcolwidth{1}{4}}
\sphinxstyletheadfamily \sphinxAtStartPar
Typ danych
\sphinxbeforeendvarwidth
\end{varwidth}%
&\begin{varwidth}[t]{\sphinxcolwidth{1}{4}}
\sphinxstyletheadfamily \sphinxAtStartPar
Klucz / Więzy
\sphinxbeforeendvarwidth
\end{varwidth}%
\\
\sphinxmidrule
\sphinxtableatstartofbodyhook\begin{varwidth}[t]{\sphinxcolwidth{1}{4}}
\sphinxAtStartPar
filmy
\sphinxbeforeendvarwidth
\end{varwidth}%
&\begin{varwidth}[t]{\sphinxcolwidth{1}{4}}
\sphinxAtStartPar
id\_filmu
\sphinxbeforeendvarwidth
\end{varwidth}%
&\begin{varwidth}[t]{\sphinxcolwidth{1}{4}}
\sphinxAtStartPar
SERIAL
\sphinxbeforeendvarwidth
\end{varwidth}%
&\begin{varwidth}[t]{\sphinxcolwidth{1}{4}}
\sphinxAtStartPar
PK
\sphinxbeforeendvarwidth
\end{varwidth}%
\\
\sphinxhline\begin{varwidth}[t]{\sphinxcolwidth{1}{4}}
\sphinxAtStartPar
filmy
\sphinxbeforeendvarwidth
\end{varwidth}%
&\begin{varwidth}[t]{\sphinxcolwidth{1}{4}}
\sphinxAtStartPar
tytul
\sphinxbeforeendvarwidth
\end{varwidth}%
&\begin{varwidth}[t]{\sphinxcolwidth{1}{4}}
\sphinxAtStartPar
VARCHAR(255)
\sphinxbeforeendvarwidth
\end{varwidth}%
&\begin{varwidth}[t]{\sphinxcolwidth{1}{4}}
\sphinxAtStartPar
NOT NULL
\sphinxbeforeendvarwidth
\end{varwidth}%
\\
\sphinxhline\begin{varwidth}[t]{\sphinxcolwidth{1}{4}}
\sphinxAtStartPar
filmy
\sphinxbeforeendvarwidth
\end{varwidth}%
&\begin{varwidth}[t]{\sphinxcolwidth{1}{4}}
\sphinxAtStartPar
rok\_produkcji
\sphinxbeforeendvarwidth
\end{varwidth}%
&\begin{varwidth}[t]{\sphinxcolwidth{1}{4}}
\sphinxAtStartPar
INT
\sphinxbeforeendvarwidth
\end{varwidth}%
&\begin{varwidth}[t]{\sphinxcolwidth{1}{4}}
\sphinxbeforeendvarwidth
\end{varwidth}%
\\
\sphinxhline\begin{varwidth}[t]{\sphinxcolwidth{1}{4}}
\sphinxAtStartPar
filmy
\sphinxbeforeendvarwidth
\end{varwidth}%
&\begin{varwidth}[t]{\sphinxcolwidth{1}{4}}
\sphinxAtStartPar
czas\_trwania
\sphinxbeforeendvarwidth
\end{varwidth}%
&\begin{varwidth}[t]{\sphinxcolwidth{1}{4}}
\sphinxAtStartPar
INT
\sphinxbeforeendvarwidth
\end{varwidth}%
&\begin{varwidth}[t]{\sphinxcolwidth{1}{4}}
\sphinxAtStartPar
NOT NULL
\sphinxbeforeendvarwidth
\end{varwidth}%
\\
\sphinxhline\begin{varwidth}[t]{\sphinxcolwidth{1}{4}}
\sphinxAtStartPar
sale
\sphinxbeforeendvarwidth
\end{varwidth}%
&\begin{varwidth}[t]{\sphinxcolwidth{1}{4}}
\sphinxAtStartPar
id\_sali
\sphinxbeforeendvarwidth
\end{varwidth}%
&\begin{varwidth}[t]{\sphinxcolwidth{1}{4}}
\sphinxAtStartPar
SERIAL
\sphinxbeforeendvarwidth
\end{varwidth}%
&\begin{varwidth}[t]{\sphinxcolwidth{1}{4}}
\sphinxAtStartPar
PK
\sphinxbeforeendvarwidth
\end{varwidth}%
\\
\sphinxhline\begin{varwidth}[t]{\sphinxcolwidth{1}{4}}
\sphinxAtStartPar
sale
\sphinxbeforeendvarwidth
\end{varwidth}%
&\begin{varwidth}[t]{\sphinxcolwidth{1}{4}}
\sphinxAtStartPar
nazwa\_sali
\sphinxbeforeendvarwidth
\end{varwidth}%
&\begin{varwidth}[t]{\sphinxcolwidth{1}{4}}
\sphinxAtStartPar
VARCHAR(50)
\sphinxbeforeendvarwidth
\end{varwidth}%
&\begin{varwidth}[t]{\sphinxcolwidth{1}{4}}
\sphinxAtStartPar
NOT NULL
\sphinxbeforeendvarwidth
\end{varwidth}%
\\
\sphinxhline\begin{varwidth}[t]{\sphinxcolwidth{1}{4}}
\sphinxAtStartPar
sale
\sphinxbeforeendvarwidth
\end{varwidth}%
&\begin{varwidth}[t]{\sphinxcolwidth{1}{4}}
\sphinxAtStartPar
pojemnosc
\sphinxbeforeendvarwidth
\end{varwidth}%
&\begin{varwidth}[t]{\sphinxcolwidth{1}{4}}
\sphinxAtStartPar
INT
\sphinxbeforeendvarwidth
\end{varwidth}%
&\begin{varwidth}[t]{\sphinxcolwidth{1}{4}}
\sphinxAtStartPar
NOT NULL
\sphinxbeforeendvarwidth
\end{varwidth}%
\\
\sphinxhline\begin{varwidth}[t]{\sphinxcolwidth{1}{4}}
\sphinxAtStartPar
seanse
\sphinxbeforeendvarwidth
\end{varwidth}%
&\begin{varwidth}[t]{\sphinxcolwidth{1}{4}}
\sphinxAtStartPar
id\_seansu
\sphinxbeforeendvarwidth
\end{varwidth}%
&\begin{varwidth}[t]{\sphinxcolwidth{1}{4}}
\sphinxAtStartPar
SERIAL
\sphinxbeforeendvarwidth
\end{varwidth}%
&\begin{varwidth}[t]{\sphinxcolwidth{1}{4}}
\sphinxAtStartPar
PK
\sphinxbeforeendvarwidth
\end{varwidth}%
\\
\sphinxhline\begin{varwidth}[t]{\sphinxcolwidth{1}{4}}
\sphinxAtStartPar
seanse
\sphinxbeforeendvarwidth
\end{varwidth}%
&\begin{varwidth}[t]{\sphinxcolwidth{1}{4}}
\sphinxAtStartPar
id\_filmu
\sphinxbeforeendvarwidth
\end{varwidth}%
&\begin{varwidth}[t]{\sphinxcolwidth{1}{4}}
\sphinxAtStartPar
INT
\sphinxbeforeendvarwidth
\end{varwidth}%
&\begin{varwidth}[t]{\sphinxcolwidth{1}{4}}
\sphinxAtStartPar
FK (filmy.id\_filmu)
\sphinxbeforeendvarwidth
\end{varwidth}%
\\
\sphinxhline\begin{varwidth}[t]{\sphinxcolwidth{1}{4}}
\sphinxAtStartPar
seanse
\sphinxbeforeendvarwidth
\end{varwidth}%
&\begin{varwidth}[t]{\sphinxcolwidth{1}{4}}
\sphinxAtStartPar
id\_sali
\sphinxbeforeendvarwidth
\end{varwidth}%
&\begin{varwidth}[t]{\sphinxcolwidth{1}{4}}
\sphinxAtStartPar
INT
\sphinxbeforeendvarwidth
\end{varwidth}%
&\begin{varwidth}[t]{\sphinxcolwidth{1}{4}}
\sphinxAtStartPar
FK (sale.id\_sali)
\sphinxbeforeendvarwidth
\end{varwidth}%
\\
\sphinxhline\begin{varwidth}[t]{\sphinxcolwidth{1}{4}}
\sphinxAtStartPar
seanse
\sphinxbeforeendvarwidth
\end{varwidth}%
&\begin{varwidth}[t]{\sphinxcolwidth{1}{4}}
\sphinxAtStartPar
czas\_rozpoczecia
\sphinxbeforeendvarwidth
\end{varwidth}%
&\begin{varwidth}[t]{\sphinxcolwidth{1}{4}}
\sphinxAtStartPar
TIMESTAMP
\sphinxbeforeendvarwidth
\end{varwidth}%
&\begin{varwidth}[t]{\sphinxcolwidth{1}{4}}
\sphinxAtStartPar
NOT NULL
\sphinxbeforeendvarwidth
\end{varwidth}%
\\
\sphinxbottomrule
\end{tabular}
\sphinxtableafterendhook\par
\sphinxattableend\end{savenotes}

\sphinxAtStartPar
\sphinxstylestrong{SQLite}


\begin{savenotes}\sphinxattablestart
\sphinxthistablewithglobalstyle
\centering
\sphinxcapstartof{table}
\sphinxthecaptionisattop
\sphinxcaption{Model fizyczny bazy danych (SQLite)}\label{\detokenize{rozdzial_3/index:id2}}
\sphinxaftertopcaption
\begin{tabular}[t]{\X{20}{100}\X{25}{100}\X{25}{100}\X{30}{100}}
\sphinxtoprule
\begin{varwidth}[t]{\sphinxcolwidth{1}{4}}
\sphinxstyletheadfamily \sphinxAtStartPar
Tabela
\sphinxbeforeendvarwidth
\end{varwidth}%
&\begin{varwidth}[t]{\sphinxcolwidth{1}{4}}
\sphinxstyletheadfamily \sphinxAtStartPar
Kolumna
\sphinxbeforeendvarwidth
\end{varwidth}%
&\begin{varwidth}[t]{\sphinxcolwidth{1}{4}}
\sphinxstyletheadfamily \sphinxAtStartPar
Typ danych
\sphinxbeforeendvarwidth
\end{varwidth}%
&\begin{varwidth}[t]{\sphinxcolwidth{1}{4}}
\sphinxstyletheadfamily \sphinxAtStartPar
Klucz / Więzy
\sphinxbeforeendvarwidth
\end{varwidth}%
\\
\sphinxmidrule
\sphinxtableatstartofbodyhook\begin{varwidth}[t]{\sphinxcolwidth{1}{4}}
\sphinxAtStartPar
filmy
\sphinxbeforeendvarwidth
\end{varwidth}%
&\begin{varwidth}[t]{\sphinxcolwidth{1}{4}}
\sphinxAtStartPar
id\_filmu
\sphinxbeforeendvarwidth
\end{varwidth}%
&\begin{varwidth}[t]{\sphinxcolwidth{1}{4}}
\sphinxAtStartPar
INTEGER
\sphinxbeforeendvarwidth
\end{varwidth}%
&\begin{varwidth}[t]{\sphinxcolwidth{1}{4}}
\sphinxAtStartPar
PK AUTOINCREMENT
\sphinxbeforeendvarwidth
\end{varwidth}%
\\
\sphinxhline\begin{varwidth}[t]{\sphinxcolwidth{1}{4}}
\sphinxAtStartPar
filmy
\sphinxbeforeendvarwidth
\end{varwidth}%
&\begin{varwidth}[t]{\sphinxcolwidth{1}{4}}
\sphinxAtStartPar
tytul
\sphinxbeforeendvarwidth
\end{varwidth}%
&\begin{varwidth}[t]{\sphinxcolwidth{1}{4}}
\sphinxAtStartPar
TEXT
\sphinxbeforeendvarwidth
\end{varwidth}%
&\begin{varwidth}[t]{\sphinxcolwidth{1}{4}}
\sphinxAtStartPar
NOT NULL
\sphinxbeforeendvarwidth
\end{varwidth}%
\\
\sphinxhline\begin{varwidth}[t]{\sphinxcolwidth{1}{4}}
\sphinxAtStartPar
filmy
\sphinxbeforeendvarwidth
\end{varwidth}%
&\begin{varwidth}[t]{\sphinxcolwidth{1}{4}}
\sphinxAtStartPar
rok\_produkcji
\sphinxbeforeendvarwidth
\end{varwidth}%
&\begin{varwidth}[t]{\sphinxcolwidth{1}{4}}
\sphinxAtStartPar
INTEGER
\sphinxbeforeendvarwidth
\end{varwidth}%
&\begin{varwidth}[t]{\sphinxcolwidth{1}{4}}
\sphinxbeforeendvarwidth
\end{varwidth}%
\\
\sphinxhline\begin{varwidth}[t]{\sphinxcolwidth{1}{4}}
\sphinxAtStartPar
filmy
\sphinxbeforeendvarwidth
\end{varwidth}%
&\begin{varwidth}[t]{\sphinxcolwidth{1}{4}}
\sphinxAtStartPar
czas\_trwania
\sphinxbeforeendvarwidth
\end{varwidth}%
&\begin{varwidth}[t]{\sphinxcolwidth{1}{4}}
\sphinxAtStartPar
INTEGER
\sphinxbeforeendvarwidth
\end{varwidth}%
&\begin{varwidth}[t]{\sphinxcolwidth{1}{4}}
\sphinxAtStartPar
NOT NULL
\sphinxbeforeendvarwidth
\end{varwidth}%
\\
\sphinxhline\begin{varwidth}[t]{\sphinxcolwidth{1}{4}}
\sphinxAtStartPar
sale
\sphinxbeforeendvarwidth
\end{varwidth}%
&\begin{varwidth}[t]{\sphinxcolwidth{1}{4}}
\sphinxAtStartPar
id\_sali
\sphinxbeforeendvarwidth
\end{varwidth}%
&\begin{varwidth}[t]{\sphinxcolwidth{1}{4}}
\sphinxAtStartPar
INTEGER
\sphinxbeforeendvarwidth
\end{varwidth}%
&\begin{varwidth}[t]{\sphinxcolwidth{1}{4}}
\sphinxAtStartPar
PK AUTOINCREMENT
\sphinxbeforeendvarwidth
\end{varwidth}%
\\
\sphinxhline\begin{varwidth}[t]{\sphinxcolwidth{1}{4}}
\sphinxAtStartPar
sale
\sphinxbeforeendvarwidth
\end{varwidth}%
&\begin{varwidth}[t]{\sphinxcolwidth{1}{4}}
\sphinxAtStartPar
nazwa\_sali
\sphinxbeforeendvarwidth
\end{varwidth}%
&\begin{varwidth}[t]{\sphinxcolwidth{1}{4}}
\sphinxAtStartPar
TEXT
\sphinxbeforeendvarwidth
\end{varwidth}%
&\begin{varwidth}[t]{\sphinxcolwidth{1}{4}}
\sphinxAtStartPar
NOT NULL
\sphinxbeforeendvarwidth
\end{varwidth}%
\\
\sphinxhline\begin{varwidth}[t]{\sphinxcolwidth{1}{4}}
\sphinxAtStartPar
sale
\sphinxbeforeendvarwidth
\end{varwidth}%
&\begin{varwidth}[t]{\sphinxcolwidth{1}{4}}
\sphinxAtStartPar
pojemnosc
\sphinxbeforeendvarwidth
\end{varwidth}%
&\begin{varwidth}[t]{\sphinxcolwidth{1}{4}}
\sphinxAtStartPar
INTEGER
\sphinxbeforeendvarwidth
\end{varwidth}%
&\begin{varwidth}[t]{\sphinxcolwidth{1}{4}}
\sphinxAtStartPar
NOT NULL
\sphinxbeforeendvarwidth
\end{varwidth}%
\\
\sphinxhline\begin{varwidth}[t]{\sphinxcolwidth{1}{4}}
\sphinxAtStartPar
seanse
\sphinxbeforeendvarwidth
\end{varwidth}%
&\begin{varwidth}[t]{\sphinxcolwidth{1}{4}}
\sphinxAtStartPar
id\_seansu
\sphinxbeforeendvarwidth
\end{varwidth}%
&\begin{varwidth}[t]{\sphinxcolwidth{1}{4}}
\sphinxAtStartPar
INTEGER
\sphinxbeforeendvarwidth
\end{varwidth}%
&\begin{varwidth}[t]{\sphinxcolwidth{1}{4}}
\sphinxAtStartPar
PK AUTOINCREMENT
\sphinxbeforeendvarwidth
\end{varwidth}%
\\
\sphinxhline\begin{varwidth}[t]{\sphinxcolwidth{1}{4}}
\sphinxAtStartPar
seanse
\sphinxbeforeendvarwidth
\end{varwidth}%
&\begin{varwidth}[t]{\sphinxcolwidth{1}{4}}
\sphinxAtStartPar
id\_filmu
\sphinxbeforeendvarwidth
\end{varwidth}%
&\begin{varwidth}[t]{\sphinxcolwidth{1}{4}}
\sphinxAtStartPar
INTEGER
\sphinxbeforeendvarwidth
\end{varwidth}%
&\begin{varwidth}[t]{\sphinxcolwidth{1}{4}}
\sphinxAtStartPar
FK (filmy.id\_filmu)
\sphinxbeforeendvarwidth
\end{varwidth}%
\\
\sphinxhline\begin{varwidth}[t]{\sphinxcolwidth{1}{4}}
\sphinxAtStartPar
seanse
\sphinxbeforeendvarwidth
\end{varwidth}%
&\begin{varwidth}[t]{\sphinxcolwidth{1}{4}}
\sphinxAtStartPar
id\_sali
\sphinxbeforeendvarwidth
\end{varwidth}%
&\begin{varwidth}[t]{\sphinxcolwidth{1}{4}}
\sphinxAtStartPar
INTEGER
\sphinxbeforeendvarwidth
\end{varwidth}%
&\begin{varwidth}[t]{\sphinxcolwidth{1}{4}}
\sphinxAtStartPar
FK (sale.id\_sali)
\sphinxbeforeendvarwidth
\end{varwidth}%
\\
\sphinxhline\begin{varwidth}[t]{\sphinxcolwidth{1}{4}}
\sphinxAtStartPar
seanse
\sphinxbeforeendvarwidth
\end{varwidth}%
&\begin{varwidth}[t]{\sphinxcolwidth{1}{4}}
\sphinxAtStartPar
czas\_rozpoczecia
\sphinxbeforeendvarwidth
\end{varwidth}%
&\begin{varwidth}[t]{\sphinxcolwidth{1}{4}}
\sphinxAtStartPar
DATETIME
\sphinxbeforeendvarwidth
\end{varwidth}%
&\begin{varwidth}[t]{\sphinxcolwidth{1}{4}}
\sphinxAtStartPar
NOT NULL
\sphinxbeforeendvarwidth
\end{varwidth}%
\\
\sphinxbottomrule
\end{tabular}
\sphinxtableafterendhook\par
\sphinxattableend\end{savenotes}

\sphinxstepscope


\chapter{Definiowanie bazy danych i wprowadzanie danych}
\label{\detokenize{rozdzial_4/index:definiowanie-bazy-danych-i-wprowadzanie-danych}}\label{\detokenize{rozdzial_4/index::doc}}
\sphinxAtStartPar
Autor: Mateusz Gałecki


\section{Definiowanie struktury bazy danych}
\label{\detokenize{rozdzial_4/index:definiowanie-struktury-bazy-danych}}
\sphinxAtStartPar
Na podstawie modeli fizycznych z rozdziału 3 przygotowałem skrypt pozwalający na utworzenie struktury tabel.

\sphinxAtStartPar
\sphinxstylestrong{PostgreSQL:}

\begin{sphinxVerbatim}[commandchars=\\\{\}]
\PYG{k}{CREATE}\PYG{+w}{ }\PYG{k}{TABLE}\PYG{+w}{ }\PYG{n}{filmy}\PYG{+w}{ }\PYG{p}{(}
\PYG{+w}{    }\PYG{n}{id\PYGZus{}filmu}\PYG{+w}{ }\PYG{n+nb}{SERIAL}\PYG{+w}{ }\PYG{k}{PRIMARY}\PYG{+w}{ }\PYG{k}{KEY}\PYG{p}{,}
\PYG{+w}{    }\PYG{n}{tytul}\PYG{+w}{ }\PYG{n+nb}{VARCHAR}\PYG{p}{(}\PYG{l+m+mi}{255}\PYG{p}{)}\PYG{+w}{ }\PYG{k}{NOT}\PYG{+w}{ }\PYG{k}{NULL}\PYG{p}{,}
\PYG{+w}{    }\PYG{n}{rok\PYGZus{}produkcji}\PYG{+w}{ }\PYG{n+nb}{INT}\PYG{p}{,}
\PYG{+w}{    }\PYG{n}{czas\PYGZus{}trwania}\PYG{+w}{ }\PYG{n+nb}{INT}\PYG{+w}{ }\PYG{k}{NOT}\PYG{+w}{ }\PYG{k}{NULL}
\PYG{p}{)}\PYG{p}{;}

\PYG{k}{CREATE}\PYG{+w}{ }\PYG{k}{TABLE}\PYG{+w}{ }\PYG{n}{sale}\PYG{+w}{ }\PYG{p}{(}
\PYG{+w}{    }\PYG{n}{id\PYGZus{}sali}\PYG{+w}{ }\PYG{n+nb}{SERIAL}\PYG{+w}{ }\PYG{k}{PRIMARY}\PYG{+w}{ }\PYG{k}{KEY}\PYG{p}{,}
\PYG{+w}{    }\PYG{n}{nazwa\PYGZus{}sali}\PYG{+w}{ }\PYG{n+nb}{VARCHAR}\PYG{p}{(}\PYG{l+m+mi}{50}\PYG{p}{)}\PYG{+w}{ }\PYG{k}{NOT}\PYG{+w}{ }\PYG{k}{NULL}\PYG{p}{,}
\PYG{+w}{    }\PYG{n}{pojemnosc}\PYG{+w}{ }\PYG{n+nb}{INT}\PYG{+w}{ }\PYG{k}{NOT}\PYG{+w}{ }\PYG{k}{NULL}
\PYG{p}{)}\PYG{p}{;}

\PYG{k}{CREATE}\PYG{+w}{ }\PYG{k}{TABLE}\PYG{+w}{ }\PYG{n}{seanse}\PYG{+w}{ }\PYG{p}{(}
\PYG{+w}{    }\PYG{n}{id\PYGZus{}seansu}\PYG{+w}{ }\PYG{n+nb}{SERIAL}\PYG{+w}{ }\PYG{k}{PRIMARY}\PYG{+w}{ }\PYG{k}{KEY}\PYG{p}{,}
\PYG{+w}{    }\PYG{n}{id\PYGZus{}filmu}\PYG{+w}{ }\PYG{n+nb}{INT}\PYG{+w}{ }\PYG{k}{REFERENCES}\PYG{+w}{ }\PYG{n}{filmy}\PYG{p}{(}\PYG{n}{id\PYGZus{}filmu}\PYG{p}{)}\PYG{p}{,}
\PYG{+w}{    }\PYG{n}{id\PYGZus{}sali}\PYG{+w}{ }\PYG{n+nb}{INT}\PYG{+w}{ }\PYG{k}{REFERENCES}\PYG{+w}{ }\PYG{n}{sale}\PYG{p}{(}\PYG{n}{id\PYGZus{}sali}\PYG{p}{)}\PYG{p}{,}
\PYG{+w}{    }\PYG{n}{czas\PYGZus{}rozpoczecia}\PYG{+w}{ }\PYG{k}{TIMESTAMP}\PYG{+w}{ }\PYG{k}{NOT}\PYG{+w}{ }\PYG{k}{NULL}
\PYG{p}{)}\PYG{p}{;}
\end{sphinxVerbatim}

\sphinxAtStartPar
\sphinxstylestrong{SQLite:}

\begin{sphinxVerbatim}[commandchars=\\\{\}]
\PYG{k}{CREATE}\PYG{+w}{ }\PYG{k}{TABLE}\PYG{+w}{ }\PYG{n}{filmy}\PYG{+w}{ }\PYG{p}{(}
\PYG{+w}{    }\PYG{n}{id\PYGZus{}filmu}\PYG{+w}{ }\PYG{n+nb}{INTEGER}\PYG{+w}{ }\PYG{k}{PRIMARY}\PYG{+w}{ }\PYG{k}{KEY}\PYG{+w}{ }\PYG{n}{AUTOINCREMENT}\PYG{p}{,}
\PYG{+w}{    }\PYG{n}{tytul}\PYG{+w}{ }\PYG{n+nb}{TEXT}\PYG{+w}{ }\PYG{k}{NOT}\PYG{+w}{ }\PYG{k}{NULL}\PYG{p}{,}
\PYG{+w}{    }\PYG{n}{rok\PYGZus{}produkcji}\PYG{+w}{ }\PYG{n+nb}{INTEGER}\PYG{p}{,}
\PYG{+w}{    }\PYG{n}{czas\PYGZus{}trwania}\PYG{+w}{ }\PYG{n+nb}{INTEGER}\PYG{+w}{ }\PYG{k}{NOT}\PYG{+w}{ }\PYG{k}{NULL}
\PYG{p}{)}\PYG{p}{;}

\PYG{k}{CREATE}\PYG{+w}{ }\PYG{k}{TABLE}\PYG{+w}{ }\PYG{n}{sale}\PYG{+w}{ }\PYG{p}{(}
\PYG{+w}{    }\PYG{n}{id\PYGZus{}sali}\PYG{+w}{ }\PYG{n+nb}{INTEGER}\PYG{+w}{ }\PYG{k}{PRIMARY}\PYG{+w}{ }\PYG{k}{KEY}\PYG{+w}{ }\PYG{n}{AUTOINCREMENT}\PYG{p}{,}
\PYG{+w}{    }\PYG{n}{nazwa\PYGZus{}sali}\PYG{+w}{ }\PYG{n+nb}{TEXT}\PYG{+w}{ }\PYG{k}{NOT}\PYG{+w}{ }\PYG{k}{NULL}\PYG{p}{,}
\PYG{+w}{    }\PYG{n}{pojemnosc}\PYG{+w}{ }\PYG{n+nb}{INTEGER}\PYG{+w}{ }\PYG{k}{NOT}\PYG{+w}{ }\PYG{k}{NULL}
\PYG{p}{)}\PYG{p}{;}

\PYG{k}{CREATE}\PYG{+w}{ }\PYG{k}{TABLE}\PYG{+w}{ }\PYG{n}{seanse}\PYG{+w}{ }\PYG{p}{(}
\PYG{+w}{    }\PYG{n}{id\PYGZus{}seansu}\PYG{+w}{ }\PYG{n+nb}{INTEGER}\PYG{+w}{ }\PYG{k}{PRIMARY}\PYG{+w}{ }\PYG{k}{KEY}\PYG{+w}{ }\PYG{n}{AUTOINCREMENT}\PYG{p}{,}
\PYG{+w}{    }\PYG{n}{id\PYGZus{}filmu}\PYG{+w}{ }\PYG{n+nb}{INTEGER}\PYG{+w}{ }\PYG{k}{REFERENCES}\PYG{+w}{ }\PYG{n}{filmy}\PYG{p}{(}\PYG{n}{id\PYGZus{}filmu}\PYG{p}{)}\PYG{p}{,}
\PYG{+w}{    }\PYG{n}{id\PYGZus{}sali}\PYG{+w}{ }\PYG{n+nb}{INTEGER}\PYG{+w}{ }\PYG{k}{REFERENCES}\PYG{+w}{ }\PYG{n}{sale}\PYG{p}{(}\PYG{n}{id\PYGZus{}sali}\PYG{p}{)}\PYG{p}{,}
\PYG{+w}{    }\PYG{n}{czas\PYGZus{}rozpoczecia}\PYG{+w}{ }\PYG{n}{DATETIME}\PYG{+w}{ }\PYG{k}{NOT}\PYG{+w}{ }\PYG{k}{NULL}
\PYG{p}{)}\PYG{p}{;}
\end{sphinxVerbatim}


\section{Wybór mechanizmów wsadowego wprowadzania danych}
\label{\detokenize{rozdzial_4/index:wybor-mechanizmow-wsadowego-wprowadzania-danych}}
\sphinxAtStartPar
Gotowe tabele musiałem uzupełnić danymi. Brałem pod uwage następujące mechanizmy wsadowego wprowadzania danych:
\begin{itemize}
\item {} 
\sphinxAtStartPar
\sphinxstylestrong{PostgreSQL:} Środowisko to oferuje bardzo dobry import danych z plików CSV za pomocą polecenia „COPY”. Jest to najszybsza metoda ale wymaga ona kontroli na formatem pliku wsadowego. Drugą opcją jest funkcja „INSERT” która pozwala na wprowadzenie wielu wierszy jednym zapytaniem.

\item {} 
\sphinxAtStartPar
\sphinxstylestrong{SQLite:} W SQlite najbardziej znaną i zalecaną formą wprowadzania danych jest instrukcja „INSERT”.

\end{itemize}

\sphinxAtStartPar
Zdecydowałem sie na stworzenie skryptu w pythonie który na podstawie prototypowych danych generuje zapytania wykorzystujące funkcje „INSERT”. Dzięki temu wygenerowany plik pasuje do baz danych


\section{Komentarz do procesu wprowadzania danych}
\label{\detokenize{rozdzial_4/index:komentarz-do-procesu-wprowadzania-danych}}\begin{enumerate}
\sphinxsetlistlabels{\arabic}{enumi}{enumii}{}{.}%
\item {} 
\sphinxAtStartPar
W pierwszej kolejności zapełnione zostały tabele „filmy” oraz „sale” aby wygenerować na nich klucze główne.

\item {} 
\sphinxAtStartPar
Później dane są wprowadzane do tabeli „seanse” aby nie wystąpiła sytuacja z naruszeniem klucza obcego.

\end{enumerate}

\sphinxstepscope


\chapter{Zapytania do bazy danych}
\label{\detokenize{rozdzial_5/index:zapytania-do-bazy-danych}}\label{\detokenize{rozdzial_5/index::doc}}
\sphinxAtStartPar
Stworzyłem funkcje w języku python obsługujące zapytania Postgres i SQLite


\section{Lokalizacja}
\label{\detokenize{rozdzial_5/index:lokalizacja}}\begin{itemize}
\item {} 
\sphinxAtStartPar
\sphinxstylestrong{Testowany Katalog:} Plik był testowany na serwerze Jupyterhub «/home/student08» oraz Postgres \sphinxhref{mailto:'student08db1/student08@postgres}{«student08db1/student08@postgres}»

\item {} 
\sphinxAtStartPar
\sphinxstylestrong{Link do repozytorium:} Plik nazywa się „zapytania\_kino.py”

\end{itemize}


\section{Opis funkcji}
\label{\detokenize{rozdzial_5/index:opis-funkcji}}
\sphinxAtStartPar
Przygotowałem po 5 zapytań dla każdej bazy danych
\begin{enumerate}
\sphinxsetlistlabels{\arabic}{enumi}{enumii}{}{.}%
\item {} 
\sphinxAtStartPar
\sphinxstylestrong{Selekcja danych i funkcje wierszowe:} Wykorzystano funkcje \sphinxcode{\sphinxupquote{UPPER()}} oraz \sphinxcode{\sphinxupquote{SUBSTR()}} w SQLite do transformacji tekstu i dat oraz\textasciigrave{}\textasciigrave{}CASE WHEN\textasciigrave{}\textasciigrave{} w PostgreSQL do sprawdzania długości filmów.

\item {} 
\sphinxAtStartPar
\sphinxstylestrong{Funkcje agregujące:} Zastosowano operacje \sphinxcode{\sphinxupquote{SUM()}}, \sphinxcode{\sphinxupquote{COUNT()}}, \sphinxcode{\sphinxupquote{MAX()}} oraz \sphinxcode{\sphinxupquote{AVG()}} połączone \sphinxcode{\sphinxupquote{GROUP BY}} oraz \sphinxcode{\sphinxupquote{HAVING}}.

\item {} 
\sphinxAtStartPar
\sphinxstylestrong{Połączenia:} Wykorzystano (\sphinxcode{\sphinxupquote{INNER JOIN}}) do relacji między trzema tabelami (filmy, sale, seanse) oraz złączenia lewostronne (\sphinxcode{\sphinxupquote{LEFT JOIN}}) wraz z funkcjami obsługi wartości pustych (\sphinxcode{\sphinxupquote{IFNULL()}} / \sphinxcode{\sphinxupquote{COALESCE()}}) do uwzględnienia rekordów pustych.

\item {} 
\sphinxAtStartPar
\sphinxstylestrong{Operatory zbiorowe:} Użyłem \sphinxcode{\sphinxupquote{UNION}}, \sphinxcode{\sphinxupquote{EXCEPT}} oraz \sphinxcode{\sphinxupquote{INTERSECT}} do filtrowania sal i repertuaru.

\item {} 
\sphinxAtStartPar
\sphinxstylestrong{Podzapytania:} Wprowadzono podzapytania skorelowane i nieskorelowane do porównywanie czasu trwania filmów ze średnią \sphinxcode{\sphinxupquote{SELECT AVG()}}.

\end{enumerate}


\section{Dokumentacja kodu}
\label{\detokenize{rozdzial_5/index:dokumentacja-kodu}}
\sphinxAtStartPar
Dokumentacja jest generowana automatycznie na podstaiwe komentarzy przez rozszerzenie Sphinx
\index{module@\spxentry{module}!zapytania\_kino@\spxentry{zapytania\_kino}}\index{zapytania\_kino@\spxentry{zapytania\_kino}!module@\spxentry{module}}\index{pg\_kategoryzacja\_filmow() (w module zapytania\_kino)@\spxentry{pg\_kategoryzacja\_filmow()}\spxextra{w module zapytania\_kino}}\phantomsection\label{\detokenize{rozdzial_5/index:module-zapytania_kino}}

\begin{fulllineitems}
\phantomsection\label{\detokenize{rozdzial_5/index:zapytania_kino.pg_kategoryzacja_filmow}}
\pysigstartsignatures
\pysiglinewithargsret
{\sphinxcode{\sphinxupquote{zapytania\_kino.}}\sphinxbfcode{\sphinxupquote{pg\_kategoryzacja\_filmow}}}
{\sphinxparam{\DUrole{n}{conn\_params}\DUrole{p}{:}\DUrole{w}{ }\DUrole{n}{dict}}}
{{ $\rightarrow$ list}}
\pysigstopsignatures
\sphinxAtStartPar
Daje każdemu filmowi kategorie powyżej i poniżej 120 minut i pokazuje w jakiej sali był seans

\end{fulllineitems}

\index{pg\_najnowsze\_filmy\_w\_malych\_salach() (w module zapytania\_kino)@\spxentry{pg\_najnowsze\_filmy\_w\_malych\_salach()}\spxextra{w module zapytania\_kino}}

\begin{fulllineitems}
\phantomsection\label{\detokenize{rozdzial_5/index:zapytania_kino.pg_najnowsze_filmy_w_malych_salach}}
\pysigstartsignatures
\pysiglinewithargsret
{\sphinxcode{\sphinxupquote{zapytania\_kino.}}\sphinxbfcode{\sphinxupquote{pg\_najnowsze\_filmy\_w\_malych\_salach}}}
{\sphinxparam{\DUrole{n}{conn\_params}\DUrole{p}{:}\DUrole{w}{ }\DUrole{n}{dict}}}
{{ $\rightarrow$ list}}
\pysigstopsignatures
\sphinxAtStartPar
Pokazuje najnowsze filmy które są grane w salach mających mniej niż 50 miejsc

\end{fulllineitems}

\index{pg\_szczegoly\_najdluzszego\_filmu() (w module zapytania\_kino)@\spxentry{pg\_szczegoly\_najdluzszego\_filmu()}\spxextra{w module zapytania\_kino}}

\begin{fulllineitems}
\phantomsection\label{\detokenize{rozdzial_5/index:zapytania_kino.pg_szczegoly_najdluzszego_filmu}}
\pysigstartsignatures
\pysiglinewithargsret
{\sphinxcode{\sphinxupquote{zapytania\_kino.}}\sphinxbfcode{\sphinxupquote{pg\_szczegoly\_najdluzszego\_filmu}}}
{\sphinxparam{\DUrole{n}{conn\_params}\DUrole{p}{:}\DUrole{w}{ }\DUrole{n}{dict}}}
{{ $\rightarrow$ list}}
\pysigstopsignatures
\sphinxAtStartPar
Pokazuje informacje o filmach którch czas trwania jest większy niż globalna średnia innych filmów

\end{fulllineitems}

\index{pg\_wspolne\_seanse\_intersect() (w module zapytania\_kino)@\spxentry{pg\_wspolne\_seanse\_intersect()}\spxextra{w module zapytania\_kino}}

\begin{fulllineitems}
\phantomsection\label{\detokenize{rozdzial_5/index:zapytania_kino.pg_wspolne_seanse_intersect}}
\pysigstartsignatures
\pysiglinewithargsret
{\sphinxcode{\sphinxupquote{zapytania\_kino.}}\sphinxbfcode{\sphinxupquote{pg\_wspolne\_seanse\_intersect}}}
{\sphinxparam{\DUrole{n}{conn\_params}\DUrole{p}{:}\DUrole{w}{ }\DUrole{n}{dict}}}
{{ $\rightarrow$ list}}
\pysigstopsignatures
\sphinxAtStartPar
Pokazuje seanse w którch film trwa ponad 150 minut oraz filmy zrobione przed 2005 rokiem

\end{fulllineitems}

\index{pg\_zaawansowane\_zestawienie\_filmow() (w module zapytania\_kino)@\spxentry{pg\_zaawansowane\_zestawienie\_filmow()}\spxextra{w module zapytania\_kino}}

\begin{fulllineitems}
\phantomsection\label{\detokenize{rozdzial_5/index:zapytania_kino.pg_zaawansowane_zestawienie_filmow}}
\pysigstartsignatures
\pysiglinewithargsret
{\sphinxcode{\sphinxupquote{zapytania\_kino.}}\sphinxbfcode{\sphinxupquote{pg\_zaawansowane\_zestawienie\_filmow}}}
{\sphinxparam{\DUrole{n}{conn\_params}\DUrole{p}{:}\DUrole{w}{ }\DUrole{n}{dict}}}
{{ $\rightarrow$ list}}
\pysigstopsignatures
\sphinxAtStartPar
Wyświetla tutuł, liczbę seansów, średnią pojemnośc sal w którch był wyświetlany oraz pokazuje tylko filmy z przynajmniej jednym seansem

\end{fulllineitems}

\index{sqlite\_filmy\_archiwalne\_i\_premium() (w module zapytania\_kino)@\spxentry{sqlite\_filmy\_archiwalne\_i\_premium()}\spxextra{w module zapytania\_kino}}

\begin{fulllineitems}
\phantomsection\label{\detokenize{rozdzial_5/index:zapytania_kino.sqlite_filmy_archiwalne_i_premium}}
\pysigstartsignatures
\pysiglinewithargsret
{\sphinxcode{\sphinxupquote{zapytania\_kino.}}\sphinxbfcode{\sphinxupquote{sqlite\_filmy\_archiwalne\_i\_premium}}}
{\sphinxparam{\DUrole{n}{db\_path}\DUrole{p}{:}\DUrole{w}{ }\DUrole{n}{str}}}
{{ $\rightarrow$ list}}
\pysigstopsignatures
\sphinxAtStartPar
Pokazuje liste filmów które został zrobione przed 2000 rokiem oraz były wyświetlane w sali o pojemności większej niż 100

\end{fulllineitems}

\index{sqlite\_pojemnosc\_sal\_dla\_dlugich\_filmow() (w module zapytania\_kino)@\spxentry{sqlite\_pojemnosc\_sal\_dla\_dlugich\_filmow()}\spxextra{w module zapytania\_kino}}

\begin{fulllineitems}
\phantomsection\label{\detokenize{rozdzial_5/index:zapytania_kino.sqlite_pojemnosc_sal_dla_dlugich_filmow}}
\pysigstartsignatures
\pysiglinewithargsret
{\sphinxcode{\sphinxupquote{zapytania\_kino.}}\sphinxbfcode{\sphinxupquote{sqlite\_pojemnosc\_sal\_dla\_dlugich\_filmow}}}
{\sphinxparam{\DUrole{n}{db\_path}\DUrole{p}{:}\DUrole{w}{ }\DUrole{n}{str}}}
{{ $\rightarrow$ list}}
\pysigstopsignatures
\sphinxAtStartPar
Pokazuje wszystkie sale w których filmy są dłuższe niż średni czas trwania wszystkich filmów

\end{fulllineitems}

\index{sqlite\_puste\_sale() (w module zapytania\_kino)@\spxentry{sqlite\_puste\_sale()}\spxextra{w module zapytania\_kino}}

\begin{fulllineitems}
\phantomsection\label{\detokenize{rozdzial_5/index:zapytania_kino.sqlite_puste_sale}}
\pysigstartsignatures
\pysiglinewithargsret
{\sphinxcode{\sphinxupquote{zapytania\_kino.}}\sphinxbfcode{\sphinxupquote{sqlite\_puste\_sale}}}
{\sphinxparam{\DUrole{n}{db\_path}\DUrole{p}{:}\DUrole{w}{ }\DUrole{n}{str}}}
{{ $\rightarrow$ list}}
\pysigstopsignatures
\sphinxAtStartPar
Pokazuje w jakich salach nie odbył się żaden seans

\end{fulllineitems}

\index{sqlite\_raport\_popularnosci\_filmow() (w module zapytania\_kino)@\spxentry{sqlite\_raport\_popularnosci\_filmow()}\spxextra{w module zapytania\_kino}}

\begin{fulllineitems}
\phantomsection\label{\detokenize{rozdzial_5/index:zapytania_kino.sqlite_raport_popularnosci_filmow}}
\pysigstartsignatures
\pysiglinewithargsret
{\sphinxcode{\sphinxupquote{zapytania\_kino.}}\sphinxbfcode{\sphinxupquote{sqlite\_raport\_popularnosci\_filmow}}}
{\sphinxparam{\DUrole{n}{db\_path}\DUrole{p}{:}\DUrole{w}{ }\DUrole{n}{str}}}
{{ $\rightarrow$ list}}
\pysigstopsignatures
\sphinxAtStartPar
Tworzy raport zawierający tytuł, liczcbę seansów i największą sale w której był pokazany

\end{fulllineitems}

\index{sqlite\_seanse\_wielkie\_litery() (w module zapytania\_kino)@\spxentry{sqlite\_seanse\_wielkie\_litery()}\spxextra{w module zapytania\_kino}}

\begin{fulllineitems}
\phantomsection\label{\detokenize{rozdzial_5/index:zapytania_kino.sqlite_seanse_wielkie_litery}}
\pysigstartsignatures
\pysiglinewithargsret
{\sphinxcode{\sphinxupquote{zapytania\_kino.}}\sphinxbfcode{\sphinxupquote{sqlite\_seanse\_wielkie\_litery}}}
{\sphinxparam{\DUrole{n}{db\_path}\DUrole{p}{:}\DUrole{w}{ }\DUrole{n}{str}}}
{{ $\rightarrow$ list}}
\pysigstopsignatures
\sphinxAtStartPar
Wyświetla wszystkie seanse po 2026 roku zmieniając tytuły na wielkie litery

\end{fulllineitems}



\renewcommand{\indexname}{Indeks modułów Pythona}
\begin{sphinxtheindex}
\let\bigletter\sphinxstyleindexlettergroup
\bigletter{z}
\item\relax\sphinxstyleindexentry{zapytania\_kino}\sphinxstyleindexpageref{rozdzial_5/index:\detokenize{module-zapytania_kino}}
\end{sphinxtheindex}

\renewcommand{\indexname}{Indeks}
\printindex
\end{document}